\documentclass[10pt,a4paper]{article}

% ----------------- Paquetes -------------------%
\usepackage[T1]{fontenc}
\usepackage[utf8]{inputenc}
\usepackage[a4paper,margin=2.5cm]{geometry}
\usepackage{mathptmx}             
\usepackage{changepage} 
\usepackage{setspace}
\usepackage[spanish]{babel}
\usepackage[labelfont=bf,labelsep=space]{caption}
\usepackage{amssymb}
\usepackage{amsmath}
\usepackage{ebgaramond}
\usepackage{placeins}
\usepackage{etoolbox}
\usepackage{setspace}
\usepackage{parskip} 
\usepackage{tcolorbox}
\usepackage{needspace}
\usepackage{xparse}
\usepackage{ocgx2}
\usepackage{hyperref}
\usepackage{hypcap}


%%%%%%%%%%%%%%%%%%%%%%%%%%%%%%%%%%%%%%%%%%%%%%%%%%%%%%%%%%%%%%%%
% --- Fija primero tus valores globales "buenos" ---
\setstretch{1.2}                 % Interlineado global
\setlength{\parskip}{0.7\baselineskip}
\setlength{\parindent}{0pt}
\newlength{\GlobalParskip}
\setlength{\GlobalParskip}{\parskip}

\newcounter{eqnum}[subsection]
\renewcommand{\theeqnum}{\arabic{eqnum}}
\newcounter{alphaitem}
\newcounter{numitem}

%% ------------------- Recuadros----------------------%%
\newcommand{\puntosclave}[1]{%
  \begin{tcolorbox}[
    colframe=black,
    title={Puntos clave},
    before upper={\setlength{\parindent}{0.5cm}\setlength{\parskip}{0.5\baselineskip}}
  ]
  #1
  \end{tcolorbox}
}

\newcommand{\modificaciones}[1]{
  \begin{tcolorbox}[
    colback=white,
    colframe=red,
    title={Modificaciones a realizar},
    before upper={\setlength{\parindent}{0.5cm}\setlength{\parskip}{0.5\baselineskip}}
  ]
  #1
  \end{tcolorbox}
}

%% ------------------- Preguntas TEST----------------------%%
% Opciones: radio (single-choice)
\NewDocumentCommand{\OptRadio}{m m m}{%
  \noindent\CheckBox[radio,name=#1,value=#2,width=1.2ex,height=1.2ex]{}%
  \;\textbf{#2.}~#3\par
}

% Opciones: checkbox (multi-choice)
\NewDocumentCommand{\OptCheck}{m m}{%
  \noindent\CheckBox[name=#1,width=1.2ex,height=1.2ex,bordercolor={0 0 0}]{}%
  \;#2\par
}

% Pregunta single-choice (radio)
% Args: 1=id, 2=enunciado, 3=opciones, 4=solución+justificación, 5=imagen (o {}), 6=origen
\NewDocumentCommand{\PreguntaTestSingle}{m m m m m m}{%
  \par\medskip
  \noindent\textbf{#1.}~#2\par\smallskip

    % Imagen (si no es {})
    \IfBlankTF{#5}{}{%
      \begin{center}
        \includegraphics[width=0.75\linewidth]{#5}
      \end{center}
    }

    \begin{Form}
      #3
    \end{Form}

    \par\smallskip
    \switchocg{sol-#1}{\fbox{Mostrar / ocultar solución}}%
    \hfill{\footnotesize #6}\par\smallskip

    \begin{ocg}{Solución}{sol-#1}{0}
      \noindent #4
    \end{ocg}
  \par\medskip
}

% Pregunta multi-choice (checkboxes)
\NewDocumentCommand{\PreguntaTestMulti}{m m m m m m}{%
  \par\medskip
  \noindent\textbf{#1.}~#2\par\smallskip

    % Imagen (si no es {})
    \IfBlankTF{#5}{}{%
      \begin{center}
        \includegraphics[width=0.75\linewidth]{#5}
      \end{center}
    }
    \begin{Form}
      #3
    \end{Form}

    \par\smallskip
    \switchocg{sol-#1}{\fbox{Mostrar / ocultar solución}}%
    \hfill{\footnotesize #6}\par\smallskip

    \begin{ocg}{Solución}{sol-#1}{0}
      \noindent #4
    \end{ocg}
  \par\medskip
}

%% ------------------- Listas----------------------%%
    % --- Lista alfabética ---
    \newenvironment{la}
      {\begin{list}{\alph{alphaitem})}{%
          \usecounter{alphaitem}%
          \setlength{\leftmargin}{1cm}%
          \setlength{\parskip}{3pt}% 
          \setlength{\itemsep}{0pt}% ← espacio entre ítems
          \setlength{\parsep}{0pt}%  ← párrafos dentro de un ítem
      }}
      {\end{list}}
      
    % --- Lista numérica ---
    \newenvironment{lnum}
      {\begin{list}{\arabic{numitem}.}{%
          \usecounter{numitem}%
          \setlength{\leftmargin}{1cm}%
          \setlength{\parskip}{3pt}% 
          \setlength{\itemsep}{0pt}% ← espacio entre ítems
          \setlength{\parsep}{0pt}%  ← párrafos dentro de un ítem
      }}
      {\end{list}}
      
% ------------------- Imágenes----------------------%
    \graphicspath{{Img/}}% Rutas por defecto 
    % --- Imagen adaptativa ---
    % Registros de dimensión definidos una sola vez
    \newdimen\imgcap
    \newdimen\imgmin
    \newdimen\imgmax
    \newdimen\imgremain
    \newdimen\imgh
    \newdimen\imgneed
    
    \newcommand{\imagenfit}[3]{%
      \addvspace{10pt}% separación antes
      \par\begingroup
        % Parámetros de composición (asignaciones, no \newdimen)
        \imgcap=1\baselineskip            % alto estimado de título+fuente
        \imgmin=0.15\textheight
        \imgmax=0.15\textheight
        \imgremain=\pagegoal
        \ifdim\pagegoal=\maxdimen \imgremain=0pt\fi
        \advance\imgremain by -\pagetotal
        \advance\imgremain by -\imgcap
        % Decisión de altura destino
        \ifdim\imgremain<\imgmin
          \newpage
          \imgh=0.20\textheight
        \else
          \imgh=\imgremain
          \ifdim\imgh>\imgmax \imgh=\imgmax \fi
        \fi
        % Reservar espacio para evitar cortes del bloque completo
        \imgneed=\imgh \advance\imgneed by \imgcap
        \Needspace{\imgneed}
        % Cuerpo
        \noindent\begin{minipage}{\textwidth}
          \centering
          \captionof{figure}{\textemdash\ #2}\par\vspace{0.3em}% título arriba
          \includegraphics[width=\linewidth,height=\imgh,keepaspectratio]{\detokenize{#1}}\par
          \vspace{0.3em}\small\textit{Fuente: }#3% fuente abajo
        \end{minipage}%
      \endgroup
      \par\addvspace{10pt}%
    }

%------------ Bloque de RECUADRO ESPECIAL -------------%
    \newenvironment{specialblock}[1]{%
      \begin{quote} % crea sangría a izquierda y derecha
      \small % texto más pequeño
      \noindent\textbf{#1}\\[-0.3em] % título en negrita
      \rule{\linewidth}{0.4pt}\\[0.6em]}{
      \end{quote}}
      
%-------------------------- Portada -----------------------%
\newcommand{\oldtitle}[1]{\def\OldTitle{#1}}
\newcommand{\changerationale}[1]{\def\ChangeWhy{#1}}

\newcommand{\changebox}{%
\begin{tcolorbox}[title={Historial de cambios del tema}]
\textbf{Título anterior:} \OldTitle\par
\textbf{Motivo del cambio:} \ChangeWhy
\end{tcolorbox}
}

%-------------------------- Author Font -----------------------%
    \newcommand{\authorfont}[1]{{\fontfamily{EBGaramond-TLF}\selectfont #1}}

%-------------------------- Anexos -----------------------%
    \makeatletter
    \newcounter{anexo}
    \renewcommand{\theanexo}{\Roman{anexo}}
    \newcommand{\anexo}[1]{%
      \clearpage
      \phantomsection
      \refstepcounter{anexo}%
      \noindent\textbf{Anexo \theanexo.- }#1\par\medskip
      \addcontentsline{stc}{section}{Anexo \theanexo.- #1}%
    }
    \makeatother

    %-------------------------- Lista de figuras (personal) -----------------------%
\makeatletter
\newcommand{\MyListOfFigures}{%
  \clearpage
  \phantomsection
  {\centering\bfseries\Large Índice de figuras\par}%
  \medskip
  % Entrada en nuestro índice de contenido (stc)
  \addcontentsline{stc}{section}{Índice de figuras}%
  % Recuperación del listado (sin cabecera propia)
  \begingroup
    \parindent=0pt\parskip=2pt
    \@starttoc{lof}%
  \endgroup
}
\makeatother

%-------------------------- Bibliografía -----------------------%
\makeatletter
% Contador para las referencias
\newcounter{myref}
% Cabecera de la bibliografía, entra en el índice 'stc'
\newcommand{\MyBibliography}{%
  \clearpage
  \phantomsection
  {\centering\bfseries\Large Bibliografía y referencias\par}%
  \medskip
  \addcontentsline{stc}{section}{Bibliografía y referencias}%
  \setcounter{myref}{0}%
}
\newcommand{\MyBibItem}[4]{%
  \refstepcounter{myref}%
  \noindent[\themyref]\ #1.\ \emph{#2}. #3. #4\par
}
\makeatother

%--------- Establecer cuatro niveles de índice --------%
    \setcounter{tocdepth}{4}   
    \setcounter{secnumdepth}{4}% numera hasta \paragraph

%%%%%%%%%%%%%%%%%%%%%%%%%%%%%%%%

\makeatletter

    %--------- Interlineado Global --------%
    \edef\GlobalBaseLineStretch{\baselinestretch}
    
    %--------- Espaciado de seccinones --------%
    \renewcommand\section{\@startsection{section}
      {1}{\z@}
      {2.5ex \@plus 1ex \@minus .2ex} % espacio antes
      {1.5ex \@plus .2ex}             % espacio después
      {\normalfont\Large\bfseries}    % formato del título
    }
    \renewcommand\subsection{\@startsection{subsection}{2}{\z@}
      {3ex \@plus 0.8ex \@minus .2ex}
      {1.5ex \@plus .2ex}
      {\normalfont\large\bfseries}
    }
    \renewcommand\subsubsection{\@startsection{subsubsection}{3}{\z@}
      {3ex \@plus 0.5ex \@minus .2ex}
      {1ex \@plus .2ex}
      {\normalfont\normalsize\bfseries}
    }
    \renewcommand\paragraph{\@startsection{paragraph}{4}{\z@}%
    {-2ex \@plus -0.5ex \@minus -.2ex}% espacio antes
    {0.8ex \@plus .2ex}% espacio después
    {\normalfont\normalsize\itshape}}% ← cursiva en lugar de negrita

    %--------- Ecuaciones --------%   
    % Variables globales para almacenar títulos y notas
    \gdef\leq@current@section{}%
    \gdef\leq@current@subsection{}%
    \gdef\leq@last@printed@section{}%
    \gdef\leq@last@printed@subsection{}%
    
    % Comando para añadir nota (invisible en el cuerpo)
    \newcommand{\eqnote}[1]{%
      \addcontentsline{leq}{leqnote}{#1}%
    }
    
    % Comando para ecuaciones
    \newcommand{\eqblock}[2]{%
      % Verificar si necesitamos imprimir el título de sección
      \ifx\leq@current@section\leq@last@printed@section\else
        \addcontentsline{leq}{leqsection}{\leq@current@section}%
        \global\let\leq@last@printed@section\leq@current@section
        \gdef\leq@last@printed@subsection{}% Reset subsección al cambiar sección
      \fi
      % Verificar si necesitamos imprimir el título de subsección
      \ifx\leq@current@subsection\leq@last@printed@subsection\else
        \ifx\leq@current@subsection\@empty\else
          \addcontentsline{leq}{leqsubsection}{\leq@current@subsection}%
          \global\let\leq@last@printed@subsection\leq@current@subsection
        \fi
      \fi
      % Ecuación
      \refstepcounter{eqnum}%
      \begin{equation*}
        #1\tag{E.\theeqnum}
      \end{equation*}%
      \addcontentsline{leq}{leqt}{[E.\theeqnum] -- #2}%
      \addcontentsline{leq}{leqm}{#1}%
    }
    
    %%% ----- Índice de ecuaciones ----------%%%
    % Tamaño para []
    \newcommand{\leqIndexFont}{\normalsize}
    \newcommand{\leqTitleFont}{\tiny}
    \newcommand{\leqNoteFont}{\tiny}
    % Cabecera de sección 
    \newcommand*\l@leqsection[2]{%
      \addvspace{25pt}% Mayor separación antes de nueva sección
      \noindent\leqIndexFont
      \makebox[\linewidth]{\rule{.38\linewidth}{0.4pt}\ \ \textbf{  \MakeUppercase{#1}}\ \ \rule{.38\linewidth}{0.4pt}}%
      \par\vspace{-10pt}
    }
    % Cabecera de subsección 
    \newcommand*\l@leqsubsection[2]{%
      \par\addvspace{15pt}%
      \noindent\leqIndexFont #1
      \par\vspace{-7pt}
    }
    % Título de ecuación 
    \newcommand*\l@leqt[2]{%
    \par\addvspace{10pt}%
      \par\noindent\leqTitleFont #1\par
    }
    % Miniatura de ecuación
    \newcommand*\l@leqm[2]{%
      \vspace{0.3\baselineskip}%
      \par\scriptsize\noindent\hspace*{1.5em}\(#1\)\par%
    }
    % Nota de ecuación 
    \newcommand*\l@leqnote[2]{%
      \vspace{0.2\baselineskip}%
      \par\noindent\leqNoteFont\hspace*{1.5em}\textbf{Nota.--} #1\par\vspace{0.5\baselineskip}%
    }
    % Generación del índice
    \newcommand{\mylistofequations}{%
      \clearpage
      \phantomsection
      % Título visual de la lista (ajústalo a tu gusto):
      {\centering\bfseries\Large Índice de ecuaciones\par}
      % Entrada en nuestro índice 'stc':
      \addcontentsline{stc}{section}{Índice de ecuaciones}%
      % Llamada a tu lista real (usa el nombre exacto de tu macro):
      \@starttoc{leq}%
    }
    
    %--------- Captura de títulos de sección --------%
    \let\leq@original@section\section
    \let\leq@original@subsection\subsection
    % Flag para saber si estamos en una sección asterisco
    \newif\ifLeqStarSection
    \LeqStarSectionfalse
    
    % Redefinir \section CON CONTROL DE ASTERISCO
    \renewcommand{\section}{%
      \@ifstar{\leq@section@star}{\@dblarg\leq@section@nostar}%
    }
    \def\leq@section@star#1{%
      \global\LeqStarSectiontrue
      \gdef\leq@current@section{#1}%
      \gdef\leq@current@subsection{}%
      \leq@original@section*{#1}%
      \global\LeqStarSectionfalse
    }
    \def\leq@section@nostar[#1]#2{%
      \global\LeqStarSectionfalse
      \gdef\leq@current@section{#2}%
      \gdef\leq@current@subsection{}%
      \leq@original@section[#1]{#2}%
    }
    % Redefinir \subsection
    \renewcommand{\subsection}{%
      \@ifstar{\leq@subsection@star}{\@dblarg\leq@subsection@nostar}%
    }
    \def\leq@subsection@star#1{%
      \global\LeqStarSectiontrue
      \gdef\leq@current@subsection{#1}%
      \leq@original@subsection*{#1}%
      \global\LeqStarSectionfalse
    }
    \def\leq@subsection@nostar[#1]#2{%
      \global\LeqStarSectionfalse
      \gdef\leq@current@subsection{#2}%
      \leq@original@subsection[#1]{#2}%
    }

    % ----- Restauración de valores Globales de estilo ----------%
    % Flags
    \newif\ifInFootnote
    \InFootnotefalse
    \pretocmd{\@footnotetext}{\InFootnotetrue}{}{}
    \apptocmd{\@footnotetext}{\InFootnotefalse}{}{}
    % Profundidad de listas (soporta anidadas)
    \newcount\MyListDepth \MyListDepth=0
    \AtBeginEnvironment{la}{\advance\MyListDepth by 1}
    \AfterEndEnvironment{la}{\advance\MyListDepth by -1}
    \AtBeginEnvironment{lnum}{\advance\MyListDepth by 1}
    \AfterEndEnvironment{lnum}{\advance\MyListDepth by -1}
    \AtBeginEnvironment{itemize}{\advance\MyListDepth by 1}
    \AfterEndEnvironment{itemize}{\advance\MyListDepth by -1}
    \AtBeginEnvironment{enumerate}{\advance\MyListDepth by 1}
    \AfterEndEnvironment{enumerate}{\advance\MyListDepth by -1}
    % Restauración diferida: hazlo al INICIO del próximo párrafo real
    \newif\if@pendingRestore
    \@pendingRestorefalse
    \newcommand{\RequestRestore}{\global\@pendingRestoretrue}
    \everypar{%
      \if@pendingRestore
        \global\@pendingRestorefalse
        \ifInFootnote\else
          \ifnum\MyListDepth=0
            \setlength{\parskip}{\GlobalParskip}%
            \setstretch{\GlobalBaseLineStretch}%
          \fi
        \fi
      \fi
    }
    % Al final de bloques:
    \AfterEndEnvironment{equation}{\RequestRestore}
    \AfterEndEnvironment{equation*}{\RequestRestore}
    \AfterEndEnvironment{la}{\RequestRestore}
    \AfterEndEnvironment{lnum}{\RequestRestore}
    % Al final de macros:
    \apptocmd{\eqblock}{\RequestRestore}{}{}
    \apptocmd{\imagenfit}{\RequestRestore}{}{}
    
\makeatother

% ===================== ToC propio con 4 niveles (único bloque) =====================
\makeatletter

% --- Escritores: añadimos una línea al .stc en cada cabecera numerada ---
% Guardamos las versiones actuales para llamar al comportamiento normal
\let\MyTOC@orig@section\section
\let\MyTOC@orig@subsection\subsection
\let\MyTOC@orig@subsubsection\subsubsection
\let\MyTOC@orig@paragraph\paragraph

% SECTION
\renewcommand{\section}{%
  \@ifstar{\MyTOC@section@star}{\@dblarg\MyTOC@section@nostar}%
}
\def\MyTOC@section@star#1{%
  \MyTOC@orig@section*{#1}% sin entrada al .stc
}
\def\MyTOC@section@nostar[#1]#2{%
  \MyTOC@orig@section[{#1}]{#2}% comportamiento normal (escribe al .toc)
  % Escribimos también nuestra línea al .stc (texto con \numberline)
  \addcontentsline{stc}{section}{\protect\numberline{\thesection}#2}%
}

% SUBSECTION
\renewcommand{\subsection}{%
  \@ifstar{\MyTOC@subsection@star}{\@dblarg\MyTOC@subsection@nostar}%
}
\def\MyTOC@subsection@star#1{%
  \MyTOC@orig@subsection*{#1}%
}
\def\MyTOC@subsection@nostar[#1]#2{%
  \MyTOC@orig@subsection[{#1}]{#2}%
  \addcontentsline{stc}{subsection}{\protect\numberline{\thesubsection}#2}%
}

% SUBSUBSECTION
\renewcommand{\subsubsection}{%
  \@ifstar{\MyTOC@subsubsection@star}{\@dblarg\MyTOC@subsubsection@nostar}%
}
\def\MyTOC@subsubsection@star#1{%
  \MyTOC@orig@subsubsection*{#1}%
}
\def\MyTOC@subsubsection@nostar[#1]#2{%
  \MyTOC@orig@subsubsection[{#1}]{#2}%
  \addcontentsline{stc}{subsubsection}{\protect\numberline{\thesubsubsection}#2}%
}

% PARAGRAPH
\renewcommand{\paragraph}{%
  \@ifstar{\MyTOC@paragraph@star}{\@dblarg\MyTOC@paragraph@nostar}%
}
\def\MyTOC@paragraph@star#1{%
  \MyTOC@orig@paragraph*{#1}%
}
\def\MyTOC@paragraph@nostar[#1]#2{%
  \MyTOC@orig@paragraph[{#1}]{#2}%
  % Si no numeras los \paragraph, cambia \theparagraph por texto plano sin \numberline
  \addcontentsline{stc}{paragraph}{\protect\numberline{\theparagraph}#2}%
}

% --- Impresor del índice propio (lee el .stc) ---
\newcommand{\MyTOC}{%
  \par\bigskip
  {\centering\bfseries\Large Índice de contenido\par}%
  \medskip
  \begingroup
    \parindent=0pt\parskip=2pt
    \@starttoc{stc}%
  \endgroup
}

\makeatother
% ===================== fin ToC propio =====================


%%%%%%%%%%%%%%


% --- Datos del tema ---
\title{Tema 3.A.7 -- La nueva economía Keynesiana\\
\large Primera y segunda generación}
\author{Marcelo Ortega}
\date{Noviembre 2025}
\newcommand{\TemaCode}{3A07}

\oldtitle{Tema 3.A.30. La nueva economía Keynesiana}
\changerationale{Misma explicación que en el Tema 3.A.6. Se espera una aproximación metodológica del expositor, no una representación sucesiva de modelos puesto que estos ya se estudian en los Temas correspondientes de Macroeconomía}


\begin{document}


\maketitle
\changebox

\newpage

% --- Página de resumen y modificaciones ---
\puntosclave{ %% Escribir puntos clave Apuntes CECO
    }
\vspace{1cm}

\modificaciones{

    NEK-1:
    \begin{lnum}
        \item MUY IMPORTANTE: Explicar rigideces nominales y rigideces reales en un mismo apartado a través del Mercado de Bienes (modelo de costes de menú, etc.) con el fin de \textbf{justifcar la intervención a través de políticas de demanda}. Explicar como interactúan las diferentes rigideces pero mantener el título del apartado como "Rigideces nominales" pues es lo más importante en este ámbito (introducir toda la argumentación de las rigideces reales como importante para amplificar los efectos de los rigideces nominales)
        \item MUY IMPORTANTE: \textbf{Justificar el desempleo involuntario} a través de las rigideces reales en el mercado de trabajo. Itroducir o priorizar los modelos del mercado de trabajo (como el modelo de salarios de eficiencia).
        \item Está bien dividido en rigideces nominales y rigideces reales, pero al principio de cada apartado es imprescindible definir que es cada rigidez y por qué los economistas de la NEK-1 las introducen: RR (desmpleo) y RN (PE de DA)
    \end{lnum}

    NEK-2:
    \begin{lnum}
        \item En la introducción definir claramente y de forma directa las diferencias entre ambas generaciones (lo que se quedan y lo que descartan) y lo que ésta última adopta del enfoque de la NMC.
        \item \textbf{IMPORTANTE:} Una de las principales críticas dirigidas hacia la fijación de precios à la CALVO es que no está microfundamentada. (Mencionar o bien en las reglas de ajuste temporal de TAYLOR, NEK-1, o en la derivación del modelo de la NEK-2);
        \item Elaborar de manera simple pero más detenidamenta las relaciones que existen entre las funciones del modelo de tres ecuaciones;
        \item Poner énfasis en el problema del zero-lower y las limitaciones que esto supuso para la creación del \textit{Quantitative easing};
        \item \textbf{RECORDAR y ENFATIZAR:} La inflación óptima en el modelo de la NEK-2 es de \textbf{cero}, esta es la única tasa de inflación que permite \textit{cerrar el círculo} y \textbf{evitar los problemas derivados de la existencia de rigideces nominales}; y,
        \item Es \textbf{MUY IMPORTANTE} incluir las características clave de los modelos de la NEK-2, en particular: el agente representativo y la medición de utilidad, que permite por tanto realizar análisis de bienestar.
    \end{lnum}

    ¿Por qué los salarios en el mercado de trabajo se representan como salarios reales? -> Por dos razones principales: 
    \begin{lnum}
        \item Porque los agentes no adolecen de ilusión monetaria;
        \item Porque la utilización de salarios nominales da lugar a salarios contracíclicos
    \end{lnum}
    }

\newpage
\MyTOC
\newpage

\mylistofequations
\clearpage

\MyListOfFigures
\clearpage


\pagenumbering{arabic}
\setcounter{page}{1}


\section{Introducción}\label{intro}
En primer lugar, podríamos preguntarnos porque es importante estudiar \textbf{X} desde un punto de vista histórico. 

Para responder a esta pregunta puede resultar muy ilustrativo el Premio Nobel de Economía de 2024, en particular, el concepto de \textit{destrucción creativa}. Un autor, o autores, puede encuadrarse en una época, resultar ser el contorno entre épocas o precisamente las ideas germinales de cambio que motivan el nacimiento de otra época, pero lo que es innegable es que ese autor es fruto del conocimiento previo que la ciencia ha alcanzado.

\imagenfit{DestruccionCreativa.png}{Ilustración: Destrucción creativa}{\href{https://www.nobelprize.org/prizes/economic-sciences/2025/press-release/}{Nobel Prize Outreach 2025. Nov 2025.}}

Si visualizamos el conocemiento económico como una escalera que se ha ido desarrollando a lo largo de su breve historia, es imprescindible conocer en que momento (peldaño) nos encontramos, es decir, cuál es la frontera del conocimiento, para poder avanzar, pero no por ello es menos importante entender y conocer el camino que nos ha traído hasta aquí, puesto que este conocimiento previo es condición necesaria para la innovación.

\imagenfit{PrescriptiveKnowledgeVSPropositionalKnowledge.png}{Ilustración: Conocimiento Prescriptivo y Conocimiento Proposicional}{\href{https://www.nobelprize.org/prizes/economic-sciences/2025/press-release/}{Nobel Prize Outreach 2025. Nov 2025.}}

\subsection{Contextualización}\label{context} 
En este Tema nos centraremos en el estudio de la macroeconomía, por ser el objeto de estudio de la Nueva Economía Keynesiana. 

\subsubsection{Relevancia}\label{importance}
La Teoría Macroeconómica es la rama de la economía que estudia los fenómenos económicos de gran escala. Es decir, aquellos que involucran a decenas de miles, cientos de millones de agentes que toman decisiones con contenido económico que se repercuten mutuamente y que a su vez inducen nuevos cambios respecto las decisiones iniciales.  Para estudiar estos fenómenos, la macroeconomía hace uso de variables agregadas, como el empleo ($L$), la inflación ($\pi$) o la renta ($Y$); y serán diferentes escuelas macroeconómicas las que tratarán de dar explicación a los fenómenos que se observan en relación a estas variables agregadas, teniendo cada una su particular punto de vista. 

Éstos fenómenos a explicar conciernen fundamentalmente los cambios en esas variables agregadas tanto en el largo plazo como en el corto plazo. Para ello, los macroeconomistas utilizan diferentes herramientas que han sido objeto de transformación gradual y que, junto con conjuntos de modelos y sus implicaciones, han dado lugar a paradigmas en el sentido de KUHN. 

\subsubsection{Historia}\label{history}
La macroeconomia, en su concepción modema, tiene su origen en la epoca de la Gran Depresión y la publicación de la \textit{Teoria General sobre la Ocupación, el lnteres y el Dinero} de John Maynard Keynes (1936). Los programas de investigación dominantes en aquel momento no daban una explicación satisfactoria a la realidad de la mayoria de las economias avanzadas en la epoca: el desempleo masivo. 

Tampoco las medidas de politica económica daban el resultado esperado. La teoría planteada por Keynes rellena este hueco, centrándose en demostrar la existencia teórica del desempleo involuntario por insuficiencia de la demanda agregada y la efectividad de las politicas activas de demanda como solución al mismo, tomando como supuesto central la falta de flexibilidad de precios y salarios. 

Si bien las ideas de KEYNES fueron acogidas con entusiasmo, existia cierta confusion sobre el mensaje central de su Teoría General y acabaría siendo la interpretación de HICKS la que calase, transformando las propuestas de KEYNES en el modelo IS-LM. 

Sin embargo, ante la contradicción de dichas teorias con la eficiencia de los mercados, a finales de los años 60's, resurgieron las críticas al enfoque de la SNC. Autores como PHELPS (1967) y FRIEDMAN (1968) concentraron sus criticas en la validacion empírica del modelo propuesto por KEYNES y sus sucesores y, concretamente, en la curva de PHILLIPS, poniendo en duda la estabilidad de la relación entre inflación y desempleo. No obstante, estos enfoques alternativos no quebraron la hegemonía del enfoque de la SNC, que se asentó de forma definitiva en los años que siguieron a la segunda Guerra Mundial.

Hubo que esperar a la Nueva Macroeconomia Clasica (NMC) para vivir una nueva transición en el campo de la macroeconomía. Liderados por LUCAS, iniciaron una revolucion cientifica con cambios radicales en los temas de investigación. la orientacion conceptual y las herramientas metodologicas. Las aportaciones de LUCAS se centraron en el estudio de los ciclos económicos, que pasaron a ser considerados resultados de decisiones de optimización de agentes con informacion imperfecta pero usada de forma óptima (HER) ante perturbaciones en la economía (LUCAS y RAPPING, 1969). 

Por tanto, se abandonan completamente las rigideces de precios y salarios incluso en el corto plazo y se plantea el vaciado continuo de los mercados. En este contexto no estaban justificadas políticas de demanda activas orientadas a la corrección de fallos de mercado. En el plano metodológico, el cambio fundamental se fragua con la llegada de la Teoría del Cicio Real (TCR), iniciada con las aportaciones de KYDLAND y PRESCOTT (1982). Siguiendo los pasos de LUCAS, la Teoría del Ciclo Real se basa en modelos microfundamentados, dinámicos (con elecciones intertemporales), en un marco de equilibrio general y estocástico (con perturbaciones reales o tecnológicas y no monetarias). 

A mitad de los anos 80's, surge una nueva corriente de pensamiento, la Nueva Economia Keynesiana (NEK), que trata de revitalizar varias de las cuestiones mas importantes del debate keynesiano y que eran completamente refutadas por la NMC:
\begin{lnum}
    \item Las rigideces de precios y salarios; y,
    \item La existencia de desempleo involuntario,
\end{lnum} 

De modo que los shocks nominates pueden tener efectos reales sobre el output y el empleo como consecuencia de fricciones que impiden el ajuste completo de los precios (no neutralidad) pero aceptando algunas de las innovaciones de LUCAS, como la necesidad de microfundamentar las decisiones de los agentes y la HER. De este modo, consiguen dar una explicación microeconómica a las rigideces planteadas por KEYNES. No existía un programa de investigación único, sino mas bien un conjunto de aportaciones variadas, que compartian una serie de caracteristicas comunes. Tal y como recogen MANKIW y ROMER en su publicación \textit{The New Keynesian Economists} (1991), los modelos neo-keynesianos comparten dos proposiciones principales: \textbf{el dinero no es neutral y los mercados no son perfectos}. Se trataba, en definitiva, de modelos de equilibrio parcial, generalmente estáticos, con mercados de competencia imperfecta (agentes no totalmente precio-aceptantes) y rigidez de precios o salarios.

\subsubsection{Problemática}\label{problem}
A lo largo de esta exposición, intentaremos dar respuesta a las siguientes preguntas:
\begin{lnum}
   \item ¿\textit{Quiénes son los autores de la NEK-1 y qué aportaciones han hecho}?
   \item ¿\textit{Quiénes son los autores de la NEK-2 y qué aportaciones han hecho}?
   \item ¿\textit{Cómo se concilian las dos generaciones de NEK y la NMC}?
   \item ¿\textit{Qué sucede ahora}? ¿\textit{Cuál es la "frontera del conocimiento"  o el state of the art}?
\end{lnum}

\subsection{Estructura de la exposición}\label{structure}
Para ello, dividiremos nuestra exposición en tres partes:
\begin{lnum}
    \item En primer lugar presentaremos las aportaciones de la NEK-1 y cuál fue su recepción;
    \item En segundo lugar, presentaremos a los autores de la NEK-2, explicaremos su linea de investigación principal y su recepción; y,
    \item Por último, presentaremos una breve recapitulación de la ideas principales y una valoración del \textit{state of the art}
\end{lnum}

\clearpage
\section{Primera generación}
\puntosclave{
    Los economistas de la NEK admiten que los modelos keynesianos previos adolecían de problemas como la falta de microfundamentación y fallos a la hora de justificar el supuesto de rigideces, aunque, al contrario que LUCAS, no consideraban estos fallos como irreparables y, por ello, los nuevos modelos neokeynesianos estaban microfundamentados.

    En el nuevo antagonismo existente entre la incipiente escuela de la NEK y la escuela de la NMC, la escuela monetarista se situaba más cerca de los neokeynesianos que de la nueva Macroeconomía Clásica.

    Apoyan la visión de la SNC, principalmente por lo que respecta a las recomendaciones de política económica (fiscal y monetaria). Este último punto es relevante en la medida en que pocos años después, la visión de los economistas neokeynesianos cambió de manera radical al abandonar totalmente la SNC y aceptar la TCR.
}
Como ya hemos introducido, los autores de la Nueva Economía Keynesiana de todo excepto homogeneos respecto a sus lineas de investigación. 

Aclaremos esto pues es crucial: los autores de la NEK-1 son autores de muchas escuelas de pensamiento diferente, algunos monetaristas, otros NMC, otros institucionalistas y/o post-keynesianos que, aceptando los fallos de la no microfundamentación de los modelos empleados en términos agregados, se disponen conciliar las aportaciones metodológicas de LUCAS (y la NMC) y la visión keynesiana respecto a las imperfecciones económicas (tanto nominales como reales). Sin embargo, cada uno centrará el estudio en casos particulares en los que se puede demostrar la existencia de rigideces como consecuencia del comportamiento optimizador de los agentes, de la existencia de fallos de mercado, de la existencia de información incompleta, etc.

Por tanto, podemos caracterizar de forma (relativamente) general a los economistas de la NEK-1 como economistas que incorporan las rigideces nominales y reales a la modelización microfundamentado de variables agregadas (macroeconómicas) desde una perspectiva de equilibrio parcial (pues se trata de modelos aislados y no un modelo unificado de comportamiento agregado como puede ser la TCR).

\subsection{Contexto histórico}

\subsubsection{Autores y Universidad de referencia}

\subsection{Método: refuerzo de la Síntesis Neoclásica}

\subsubsection{Lado de la oferta}

\subsection{Rigideces nominales}
Uno de los asuntos fundamentales de los primeros trabajos de la NEK era la justificación de las rigideces nominales en precios, que, a su vez, era uno de los temas centrales de la teoría keynesiana. 

\textbf{Definición.-} Las rigideces nominales de los precios a nivel agregado surgen cuando las empresas no ajustan sus precios ante variaciones generalizadas de la demanda, y en cambio ajustan sus cantidades. 

Las rigideces fueron generalmente refutadas por el marco neoclásico bajo la premisa de que los mercados estaban en equilibrio y funcionaba en competencia perfecta (aunque autores como PIGOU analizaron profundamente los efectos de fricciones reales sobre la economía, no tanto las nominales). Con la publicación de KEYNES, las rigideces (fricciones nominales) pasaron a ser un elemento central en la explicación del desempleo, desde ahora visto como un fenómeno involuntario; también así el marco de la SNC.  Estas pequeñas fricciones que dificultan el ajuste son, por ejemplo, los costes de menu, o costes asociados a cualquier variación de los precios de una empresa, desde la impresión de nuevas listas de precios (o menus, en el caso de restaurantes) hasta el miedo a perder clientes, a iniciar una guerra de precios, etc. 

El hallazgo fundamental es que fricciones como las anteriores, que para una empresa suponen pequeñas perdidas perfectamente asumibles, pueden tener un efecto importante a nivel agregado, porque cada empresa individualmente tiene pocos incentivos para ajustarse inmediata y completamente ante pequeñas perturbaciones haciendo que, a nivel agregado, esta ausencia de ajustes individuales genere una rigidez del nivel general de precios. Por tanto, los efectos agregados de estas fricciones, que individualmente pueden tener poca importancia, son una considerable rigidez nominal del indice general de precios y la consiguiente posibilidad de fluctuaciones económicas tras una perturbación de demanda. De hecho, ROMER (1993) explica por que fue tan lento el progreso en investigar la no neutralidad de las perturbaciones de demanda agregada a raíz de este nuevo enfoque: se empezó por buscar la explicación en grandes desviaciones del modelo walrasiano, en vez de buscar grandes efectos de pequeñas desviaciones. 

Hay dos vías alternativas/complementarias para explicar estas fricciones nominales: la racionalidad incompleta de los agentes y los costes de menú. 

\subsubsection{Modelo de Costes de Menú: MANKIW}
Los costes de menú son un concepto amplio que engloba todos los posibles costes en que puede incurrir una empresa cuando altera el precio de venta de su producto, aunque generalmente estos costes son pequeños. 
Pueden ser administrativos, como el coste físico  de imprimir nuevas listas de precios y distribuirlas; o también de información, como decidir cual es el nuevo precio óptimo, para lo cual tiene que recoger información, procesarla, analizarla y llegar a una decision. Los costes de menu tambien incluyen costes no administrativos como por ejemplo el riesgo de perder una parte de los clientes habituales, que contrariados por frecuentes variaciones del precio deciden buscar otro proveedor mas estable, el riesgo de comenzar una guerra de precios o el coste de la publicidad. 

Estos coste de menú suelen representarse como una constante exógena ($\theta$) y pueden ser mayores o menores que la pérdida que enfrentan si no cambian su precio. 

Si estos son mayores, decidirá mantener su precio, realizando el ajuste en la cantidad, lo que genera fluctuaciones económicas y con ello variaciones del bienestar social (positivas o negativas) mucho mayores que la pérdida que sufre la empresa. Ademas, este resultado se obtiene con agentes racionales, puesto que la empresa busca maximizar su beneficio pero esta sujeta a unos costes de ajuste que restringen su campo de maniobra. El problema que enfrenta la empresa puede representarse analíticamente como:

\eqblock{
\begin{aligned}
    &\text{En $t_0$, la empresa ($i$) planifica:}\;\;q_{i}^e(p_{i})=E[q_{i}(p_i)\;|\;\Omega_t]\\[1em]
    &\Longrightarrow\boxed{\max_{\pi}{\pi_i=p_i\;q_{i}^e-c\;q_{i}^e}\;;\quad \text{s.a.}\; \pi_i\ge0}
\end{aligned}
}{Problema de maximización de la empresa. Fricciones nominales: Modelo Costes de menú}

Es decir, la empresa estima (HER) cuál y cómo será la curva de demanda que enfrentará en el periodo siguiente y, en base a ello define su política de precios. Durante el periodo de ejecución operativa, la demanda observada de la empresa puede haber experimentado un shock no previsto y, por tanto, enfrentar esta una decisión de:
\begin{la}
    \item Cambiar su política de precios durante el periodo de ejecución asumiendo los costes de menú; o,
    \item Llevar a cabo un ajuste de cantidad y no modificar su precio.
\end{la}

La decisión que tome dependerá exclusivamente de la magnitud de costes de menú que enfrente ya que esta tomará una decisión u otra si los costes de menú son superiores (inferiores) a la pérdida (ganancia) de beneficios que prevé ante este shock no esperado de la demanda. Gráficamente, la pérdida de beneficios se puede representar como el triangulo sombreado.\footnote{Como vemos, una condición necesaria para poder llevar a cabo esta decisión alternativa es la existencia de un cierto grado de poder de mercado (fricciones reales), aunque entraremos más en detalle en el siguiente apartado.
}

\imagenfit{CostesMenu.png}{Costes de menú}{Apuntes ICEX-CECO. Tema 3.A.7}

Si los precios permanecen estáticos (rigidez nominal), la pérdida social será mayor que la pérdida individual. 

\subsubsection{Modelo de Racionalidad Incompleta: AKERLOFF-YELLEN} 
Otra aportación muy relevante en el campo de las rigideces nominales agregadas es el presentado por AKERLOF y YELLEN (1985), que, con la intención de desvincularse de la necesidad de definir unos costes de menú cuantificables y el alcance de una decisión subóptima como consecuencia de la existencia de unos costes, generalizan este comportamiento mediante la introducción del concepto de racionalidad incompleta o cuasirracionalidad. 

\textbf{Definición.-} El comportamiento cuasirracional es aquel en que un agente no altera sus decisiones ante una perturbación si la pérdida prevista en su función objetivo es de segundo orden (pequeña). 

Si suponemos que la empresa se encuentra inicialmente maximizando su beneficio (curva $A$) a un precio $D$. Ante una perturbación adversa de la demanda, se desplaza la función de beneficio (a la curva $A'$) por lo que hay un nuevo precio que maximizaría el beneficio es $C$. Pero si la empresa no altera el precio y lo mantiene en $D$ (el mismo precio que antes de la perturbación), la reducción del beneficio ($B_1-B_2$) puede considerarse de segundo orden por lo que podría decidir mantenerse con este.

\imagenfit{RacionalidadIncompleta.png}{Racionalidad Incompleta. AKERLOFF-YELLEN}{Apuntes ICEX-CECO. Tema 3.A.7}

AKERLOF y YELLEN (1985) aplican este concepto a un modelo macroeconómico de competencia monopolística. Suponen que una proporción de las empresas son cuasiracionales y ante una perturbación monetaria no se ajustan. Los autores calculan que la diferencia entre su beneficio y el de las empresas continuamente maximizadoras es una función de segundo orden de la perturbación, pero este comportamiento provoca una variación signiticativa en las variables reales agregadas (empleo y producción, pérdida de primer orden). Por tanto, si una determinada proporción de empresas no ajustan sus precios ante un shock de demanda, aunque a nivel individual la función de pérdidas es de segundo orden, a nivel agregado pueden ser de primer orden.\footnote{La critica a la argumentación de los efectos de la cuasirracionalidad es que el modelo se cumple en una situación estática donde todas las empresas están inicialmente en su precio óptimo. En un modelo dinámico y estocástico, los agentes que no sigan reglas optimizadoras no van a estar casi nunca en una situación de equilibrio, por tanto, es posible que los costes que sufren no sean pequeños. No obstante, AKERLOF y YELLEN (1991) desarrollan un modelo para refutar esta crítica. Según los autores, las pérdidas son pequeñas también en un entorno dinámico y estocástico y, por tanto, no se pueden despreciar posibles conductas que no sean completamente maximizadoras a corto plazo (si deben serlo a largo plazo). Esta argumentación defiende la consistencia teórica de conductas propuestas en modelos keynesianos y que fueron rechazadas por no ser completamente maximizadoras.
}

\subsubsection{Modelo combinado de efectos agregados. BLANCHARD-KIYOTAKI}
Por su parte, BLANCHARD y KIYOTAKI (1987) formalizan de una manera más general los efectos macroeconómicos que pueden tener estos efectos (comportamiento cuasirracional y costes de menú).

Para ello, desarrollan un modelo basado en competencia imperfecta y consideran que los costes privados de la rigidez nominal son distintos de los costes sociales porque la rigidez nominal supone una externalidad de la demanda agregada. El beneficio social sería mayor si todas las empresas disminuyeran su precio ante un shock negativo de la demanda, pero no existen incentivos a hacerlo. El modelo plantea que la demanda percibida por la empresa $i$ se puede expresar como:

\eqblock{
\begin{aligned}
    &\text{Partiendo de la ecuación cuantitativa del dinero:}\; M^d\;V=P\;Y^d;\\[1em]
    &\text{(Consideramos $V$ constante):}\quad \boxed{Y_i^d=\underbrace{\left(\frac{M^d}{P}\right)}_{\text{\scriptsize Demanda Agregada}}\;\cdot\;\underbrace{\left(\frac{P_i}{P}\right)^{-\theta}}_{\substack{\text{\scriptsize Precio relativo y}\\ \text{\scriptsize sustituibilidad $\theta$}}}}\\[1em]
    &\text{Entonces, por ejemplo:}\\[1em]
    &\downarrow\text{\scriptsize Transacciones}\to\;\;\downarrow M^d\overbrace{\Longrightarrow}^{{\text{\scriptsize Descenso de DA}}} \downarrow P_i\;\;\text{pero,}\;\not\downarrow\;:\underbrace{\{P_i^t\sim P_i^{t-1}\}_{\perp}}_{{\substack{\text{\scriptsize Coste de $2^{\text{o}}$ orden}\\\text{\scriptsize Comportamiento cuasirracional}}}}\overbrace{\Longrightarrow }^{\text{\scriptsize Externalidad}} \underbrace{P_{\perp}}_{\text{\scriptsize Coste de $1^{\text{r}}$ orden}}
\end{aligned}
}{Demanda agregada de la empresa $i$ en función de la DA general y del precio relativo. Fricciones nominales: BLANCHARD-KIYOTAKI}

Por lo que cambios en la demanda agregada\footnote{En realidad, si somos rigurosos sería en el stock real de dinero; pero podemos utilizarlo de forma equivalente si consideramos un contexto de equilibrio general
} desplazan la curva de demanda de la empresa:
Por ejemplo, si disminuye la demanda agregada (disminuye la demanda de dinero, por ejemplo, por diminución de las trasacciones) y la empresa no ajusta el precio la empresa incurre en un coste de segundo orden, ya que el precio relativo no se ajusta al nuevo nivel de precios ($P$), pero la externalidad de la demanda hace que la rigidez de precios de la empresa ($P_i$) contribuya a la rigidez de precios del conjunto de la economia ($P$), por aplicación análoga de este mismo razonamiento a los precios del resto de empresas ($P_{-i}$). Por tanto, la externalidad surge porque el ajuste del precio individual contribuye al ajuste del nivel agregado de precios (ya que el precio de la empresa forma parte del índice de precios) y, como cada empresa ignora los beneficios agregados del ajuste de los precios y se enfrenta, únicamente, a unos costes de segundo orden, esto se traslada al conjunto de la economía de forma que se experimenta una pérdida (costes) de primer orden en el Saldo Real ($M/P$, DA). que afecta a todas las empresas.

\subsubsection{Reglas de ajuste de los precios}
Si existen restricciones lo suficientemente importantes como para justiticar que las empresas no alteren sus precios, una crítica directa que podríamos dirigir sería que en cualquiera de los modelos presentados se debería considerar desde una perspectiva dinámica, pues el comportamiento racional de los agentes les llevaría en cualquier caso a considerar su maximización intertemporal y, como los costes solo se ejecutan una vez (ante una variación de precio) podría ser rentable asumirlos. 

Ahora bien, aunque existe multitud de literatura (por ejemplo AKERLOOF y YELLEN) que justifican la rigidez nominal también desde un punto de vista dinámico, lo que nos interesa de esta crítica es una pregunta implícita, pues realmente sabemos que (o generalmente aceptamos) que los precios a nivel particular son (relativamente) rígidos pero no inamovibles, de hecho, una Botella de Agua AQUAREL costaba hace diez años 0,33€, y ahora está alrededor de 0,59€, ¿a qué se debe este incremento de precio? ¿cuándo y por qué cambian los precios las empresas?\footnote{\href{https://www.huffingtonpost.es/virales/hace-comparativa-valian-productos-super-10nos-reflexionar.html}{\textit{Hace una comparativa de cuánto valían los productos de un super hace 10 años y ahora: para reflexionar}. Huffington Post (2024)}} Esta crítica resultó ser profundamente constructiva, pues abrió un nuevo campo de investigación, \textit{¿cómo modifican sus precios las empresas?} De hecho, identificamos en la exposición dos preguntas clave:
\begin{lnum}
    \item ¿Cómo y con qué frecuencia se ajustan los precios? 
    \item ¿Qué reglas siguen las empresas para ajustar sus precios?
\end{lnum} 

\paragraph{Reglas de ajuste temporal: FISCHER y TAYLOR} 
Las reglas de tiempo consideran que los precios se modifican en función de la evolución de los acontecimientos. 
Estas reglas plantean dos cuestiones a resolver:
\begin{lnum}
    \item ¿Cuál es el periodo de tiempo entre decisiones? y,
    \item ¿Todas las empresas siguen la misma periodicidad?
\end{lnum} 

\textsc{Periodicidad: Solapamiento o sincronización}

Respecto a la segunda cuestión, lo principal será determinar si:
\begin{lnum}
    \item \textit{Solapamiento}. Si existe heterogeneidad de empresas, no todas determinan los precios en el mismo momento por lo que existe solapamiento. Esto implica que el efecto real perdura mucho mas tiempo puesto que ninguna empresa se ajusta totalmente desde el principio por miedo a cambiar demasiado su precio relativo.\footnote{Intuitivamente el solapamiento de los intervalos parece inevitable en una economía formada por cientos de miles de empresas con distinta actividad y sometidas por tanto a muy distintas perturbaciones y restricciones. La heterogeneidad parece motivo suficiente y ello es lógico que en modelos con empresas homogeneas sometidas a las mismas perturbaciones se obtenga que el único equilibrio estable es aquel en que todas las empresas han sincronizado sus intervalos (BALL, 1987).
    } Si suponemos que las empresas revisan sus precios cada dos periodos, podemos establecer que la mitad de éstas lo revisarán en $t_{i}$ y la otra mitad en $t_{i+1}$, por tanto, ante un shock nominal en $t_i$, solo una porción ajustará sus precios por lo que el shock nominal se traslada a las variables reales de una manera más prolongada en el tiempo; por el contrario,
    \item \textit{Sincronización}. Si existe sincronización el efecto real de una perturbación solo dura el tiempo que se tarda en tomar la próxima decisión de variar precios, y en ese momento, todos los precios se ajustan anulando el efecto real. De modo que los costes de menu no explicarían los efectos reales de un shock negativo sobre la economía.
\end{lnum}

Por lo tanto, seria necesario introducir explícitamente algun tipo de heterogeneidad entre las empresas para que el solapamiento se justifique formalmente. La forma mas directa es introducir perturbaciones especificas o idiosincráticas de suficiente entidad como para obligar a la empresa que la sufre a ajustar su precio en el momento y no esperar a las demas. El supuesto anterior hace referencia a las reglas de estado en vez de reglas temporales. Es interesante citar el análisis realizado por BALL y CECCHETTI (1988), en el que destacan como causas del solapamiento el hecho de que las empresas analizan las ventajas de recopilar información antes de decidir sobre su precio, por lo que tienen motivos para esperar a que sus competidoras fijen sus precios y romper la sincronización. De este modo, el ajuste a nivel agregado puede ser lento aunque los precios cambien con frecuencia. 

Respecto a los efectos del solapamiento, los trabajos de FISCHER (1977) y TAYLOR (1980) indican que el solapamiento de los salarios nominales permite que el efecto de las perturbaciones nominales sea real y duradero. Lo mismo se obtiene si hay solapamiento en la fijación de precios, puesto que lo importante es que haya rigidez nominal duradera de algún tipo. La causa de este resultado es que si las empresas cambian su precio a intervalos fijos (un mes) y una empresa tras otra, cuando a una empresa le llega el momento de cambiar su precio no va a ajustarlo completamente a la perturbación porque eso desviaría mucho su precio relativo del inicial, perjudicando su competitividad, de modo que lo modificará solo en parte. Esto ocurrirá con la siguiente empresa y así sucesivamente. El resultado es que se tarda mucho tiempo en ajustar completamente todos los precios a la perturbación, y durante ese tiempo se producen efectos reales que van desapareciendo al mismo ritmo al que los precios se ajustan.

\textsc{Periodicidad}

Respecto al intervalo de tiempo entre decisiones, en general se ha considerado que lo mas sencillo es utilizar un intervalo constante, suponiendo que la empresa se reune a intervalos fijos para revisar su precio. La longitud del intervalo se puede obtener endógenamente comparando los costes marginales (los costes de menú) y beneficios marginales (disminución de pérdidas) de reducir el intervalo en una unidad de tiempo.\footnote{Asi lo hacen BALL, MANKIW y ROMER (1988), resultando que el intervalo es una función inversa de la tasa de inflación, de la varianza de la demanda agregada y de la varianza de las perturbaciones particulares de la empresa. De aqui se obtiene la implicación de que la curva de Phillips a corto plazo tendrá mayor pendiente a medida que aumenta la inflación, puesto que la velocidad de ajuste de los precios aumenta, y las variaciones de la demanda tienen menor efecto real. 

Además, al prolongarse en el tiempo, las empresas no tendrán incentivos a llevar a cabo un ajuste total del precio, sino que lo harán de manera gradual para no perder clientes, por tanto el efecto real podrá perpetuarse incluso más de dos periodos; tiempo que, bajo la premisa de solapamiento de dos periodos, hemos establecido anteriormente.
}

\paragraph{Reglas de ajuste situacional (o estado-dependiente)}
Aparte de las reglas de ajuste según el tiempo, otros trabajos consideran que a la hora de fijar los precios son más importantes las reglas que dependen del estado de la naturaleza. En este caso, no hay un periodo de tiempo definido para cambiar los precios, sino que cambian cuando la situación rebasa ciertos umbrales o bandas. No realiza, por tanto, ajustes pequeños del precio.\footnote{Los modelos de CAPLIN y SPULBER (1987) o SHESHINSKI y WEISS (1977) responden a esta logica. La decision de cambiar los precios es endógena y en cada momento del tiempo los agentes deciden variar el precio de acuerdo con un análisis coste beneficio. Por tanto, este tipo de modelos predicen que la probabilidad de un cambio en el precio cambia según el estado de la economía. Desde el punto de vista de la racionalidad económica son interesantes ya que asumen que los agentes basan sus decisiones en un análisis coste beneficio. También es posible e intuitivamente atractivo un modelo que combine reglas dependientes del tiempo y de la situación (reglas mixtas). La empresa revisa el precio con una determinada regularidad pero si durante el periodo entre ajustes se produce una perturbación significativa en términos de una regla de estado, entonces la empresa ajusta el precio sufriendo el coste de menú.
}

\subsection{Rigideces reales}
Estas fricciones nominales que acabamos de presentar hay que entenderlas en un contexto de competencia imperfecta (fricciones reales), que introduce un amplio abanico de posibilidades para explicar la falta de respuesta de los precios ante variaciones de la producción. Esto es un elemento clave para los economistas de la NEK, pues entienden que las fricciones nominales no son suficientes para causar efectos reales significativos, sino que se requiere de rigideces reales.\footnote{BALL y ROMER (1990) demostraron que rigideces nominales sustanciales pueden provenir de una combinación de rigideces reales y pequeñas fricciones nominales. De hecho, MANKIW y ROMER (1991) identifican la interacción entre las imperfecciones nominales y reales como uno de los hechos distintivos de la NEK.
} De la misma forma, si todos los precios de una economia fueran completamente flexibles, un shock nominal dejaría el equilibrio real de la economia sin cambios. 

Por tanto, la rigidez real no implica una rigidez nominal y, de la misma forma, sin una fuente independiente de rigideces nominales, los precios se ajustan completamente a los shocks nominales a pesar de la existencia de rigideces reales. Ahora bien, las rigideces reales en el mercado de bienes y en el mercado de trabajo magnificarán la no neutralidad que resulta de las pequeñas fricciones nominales. 

ROMER (1993) se centra en las conclusiones extraídas del gráfico de los costes de menú. Para que un shock tenga consecuencias importantes no basta con incorporar pequeñas fricciones nominales en un contexto de competencia imperfecta, sino que es necesario que existan rigideces reales. Dado que los costes marginales caen mucho en las recesiones en condiciones normales, las empresas tendrán generalmente incentivos a modificar sus precios, por lo que para que los precios no se modifiquen tienen que confluir algunos de los dos factores: que los costes marginales no caigan porque exista una fuente de rigidez real que lo impida o que la elasticidad de la demanda caiga mucho ante un shock negativo.

\paragraph{Mercado de bienes}
En el mercado de bienes, las imperfecciones reales son fundamentalmente consecuencia del poder de mercado. El objetivo fundamental será entonces determinar por qué las pendientes de Ingresos Marginales y Costes Marginales son relativamente planas o, más en asbtracto, cómo interactúan y cuánto maginifican las rigideces reales el efecto de las rigideces nominales, y viceversa. 

Por ejemplo, volviendo al modelo de los costes de menú ilustrado en el apartado anterior, sabemos que el triángulo sombreado requiere la existencia de un cierto grado de poder de mercado (competencia imperfecta, es condición necesaria) pero, además, el area será una función de las pendientes del coste marginal y del ingreso marginal. De la misma forma, estas están determinadas por los rendimientos a escala y el poder de mercado de la empresa (respectivamente):
\begin{la}
    \item \textit{Costes marginales}. La tecnología (rendimientos) determinará la pendiente de la curva de costes marginales, si la empresa debe ajustarse a un shock de demanda (negativo), estaremos frente a un escenario de exceso de capacidad instalada, situándonos en la zona decreciente de la función de $CMg_i$ (rendimientos mixtos, curva de costes marginales en forma de $U$). Por tanto, experimentaremos una traslación de la curva de costes marginales sectorial, encareciendo la producción;
    \item \textit{Ingresos marginales}. La pendiente de la curva se ve directamente afectada por el poder de mercado de la empresa: cuánto mayor sea el poder de mercado de la empresa, menos elastica será la demanda y menos incentivo se tiene a modificar el precio. Otra forma de analizar el efecto sobre el ingreso marginal es ver cuánto se desplaza la curva del ingreso marginal cuando cae la demanda. Este desplazamiento va a depender de la sensibilidad de la elasticidad de la demanda. Si la elasticidad de la demanda cae ante una bajada de la producción, entonces mayor será el desplazamiento de la curva de ingresos marginales y menor será el incentivo a modificar los precios.
\end{la}

\textsc{Por el lado de los costes}

Por el lado de los costes, si suponemos que existe un cierto grado de fricciones reales (presencia de competencia imperfecta, rendimientos mixtos, etc.), podemos ilustrar gráficamente como éstas interactúan y amplifican el efecto de las fricciones nominales exploradas antes:

\imagenfit{RigidecesReales_Bienes.png}{Economías de aglomeración. DIAMOND, 1993}{Apuntes Víctor Begué. Tema 3.A.30}

Este simple modelo nos muestra que, si existe competencia imperfecta en el mercado, ante una variación (shock negativo) de la curva de demanda, la curva de costes marginales se desplazará hacia arriba, por ejemplo por una curva de costes medios en forma de U. De la misma forma, este fenómeno muestra la posibilidad de que una economía tenga equilibrios de baja y alta actividad, que es posible si el mercado es de competencia imperfecta, pues si no existe subastador walrasiano, no todos los agentes se pueden poner de acuerdo al mismo tiempo a la hora de intercambiar los bienes, no alcanzándose el equilibrio walrasiano. En las épocas de alta actividad, se compran los inputs de manera menos costosa que en los de baja actividad por lo que los costes medios caen, esto desplaza el coste marginal hacia abajo en los periodos expansivos por lo que el precio no aumenta. Existen otros modelos, complementarios o alternativos, que permiten explicar la existencia de rigideces reales por el lado de los costes en el mercado de bienes, entre otras:
\begin{lnum}
    \item \textit{Las barreras de entrada}. La empresa consigue mantenerse en el tramo decreciente de sus costes. Suele consistir en que la empresa impone un precio de venta que previene la entrada de nuevas empresas, haciendo nulos los beneficios de la potencial entrante. Esto hace, generalmente, que la empresa ya instalada se sitúe en el tramo decreciente de su curva de costes totales medios, tratando al tiempo de mantener o aumentar su producción. Puede hacerlo si tiene por ejemplo un exceso de capacidad instalada;
    \item \textit{Los plazos de entrega variable}. Supone que la empresa puede modificar el plazo de entrega de su mercancia, alargándolo, para evitar aumentar el precio del producto ante un aumento de la demanda, evitando un aumento del coste marginal. Para ello precisa una relación alta entre el precio del bien y el plazo de entrega (con elasticidad demanda-precio alta y elasticidad de demanda-entrega baja).\\
    Esta teoría es, además, muy importante ya que permite conciliar la rigidez observada con un comportamiento competitivo de las empresas, pues en este caso, es la disposición del cliente a esperar lo que permite la estabilidad de precios.
\end{lnum}

\textsc{Por el lado de los ingresos}

Un ejemplo muy sencillo para ilustrar como afectan las rigidecs reales a la cuantificación de la pérdida (y la amplificación de las rigideces nominales) es a través de la consideración de un mercado monopolístico en el que el monopolista fija el precio en base a su poder de mercado:

\eqblock{
\begin{aligned}
    &\text{En el corto plazo, $\bar K$, por tanto, el factor variable es $L$,}\\[1em]
    &P=\underbrace{\frac{w}{PMg_L}}_{\text{\scriptsize $CMg$}}\;\cdot\;\underbrace{\frac{1}{1+\varepsilon_p^{-1}}}_{\text{\scriptsize Índice de Lerner}}=\boxed{\frac{w}{PMg_L}\;\cdot\;(1+\mu)}\\[0.5em] 
    &\text{donde, $\mu\equiv$ mark-up sobre el precio}
\end{aligned}
}{Fricciones reales y elasticidad de la demanda.}

Este sería un caso en el que la elasticidad precio de la demanda juega un factor fundamental en la magnitud de las rigideces reales. A mayor producción, mayor elasticidad precio de la demanda, siendo la cantidad demandada más sensible a variaciones en el precio. Si esto es así, la empresa tendría problemas para aumentar el precio a medida que aumenta la producción, pues a mayor elasticidad precio de la demanda, mayor pérdida (reducción de la cantidad demandada) por el aumento del precio. 

Igual que desde la perspectiva de los costes, existen otros modelos, complementarios y/o alternativos, que explican la interacción entre rigideces nominales y rigideces reales, entre otras:
\begin{lnum}
    \item \textit{Los mercados de clientes}. Es posible que en muchos sectores se desarrolle una relación de confianza entre productor y consumidor, debido a los costes que supone para el consumidor buscar el proveedor mas ventajoso para cada compra. Asi, ante pequeños aumentos de precio, los consumidores no cambian de proveedor; y,
    \item \textit{La colusión contracíclica}. Si el mercado es oligopolístico, el grado de colusión es probable que sea contracíclico, pues las empresas competirán más en etapas de expansión para ganar cuotas de mercado, reduciendo el margen y no el precio, de manera que, ante un aumento de la demanda, las empresas querrán aumentar su cuota de mercado y, a la hora de competir, en lugar de modificar el precio, reduciran su margen y aumentarán la competencia.
\end{lnum}

\paragraph{Mercado de crédito}
Por otro lado, también podemos estudiar como la existencia de rigideces reales en el mercado de crédito afectan a las fricciones nominales, y viceversa. Podemos modelar las fricciones reales en el mercado de crédito como consecuencia de la presencia de información asimétrica, que da lugar a un equlibrio con racionamiento de crédito (sub-óptimo). El modelo canónico en este campo es el Modelo de STIGLITZ-WEISS (1981).\footnote{Se puede ver el desarrollo analítico completo de este modelo en [\authorfont{Apuntes ICEX-CECO Tema 3.A.7}]} 

De forma simplificada podemos considerar que existen dos agentes en el mercado de crédito:
\begin{la}
    \item El \textit{prestamista}. El banco ofrece crédito a cambio de la devolución de su capital e intereses o, en caso de insolvencia, de la ejecución de las garantías reales que se hayan depositado para solicitar el préstamo. Sin embargo, consideraremos que siempre será de mayor interés para el banco la recuperación del préstamo mediante el pago acordado, ya que la ejecución de garantías le supone un coste y no llega a cubrir los gastos en los que ha incurrido al otorgar el préstamo; y, por otro lado,
    \item El \textit{prestatario}. Los prestatarios pueden clasificarse en dos tipos (en función de su riesgo):
    \begin{la}
        \item \textit{Alto riesgo}. Los prestatarios de alto riesgo son aquellos que acometen proyectos arriesgados y prometen una alta rentabilidad; y,
        \item \textit{Bajo riesgo}. Los prestatarios de bajo riesgo son aquellos que no esperan grandes fluctiaciones en sus flujos de caja y, por tanto, están dispuestos a solicitar créditos para mejorar su rendimiento siempre y cuando sean plenamente conscientes de que pueden hacer frente al pago.
    \end{la}
\end{la}

El banco no será capaz de distinguir entre estos dos tipos de prestatario, por tanto, se verá obligado a fijar un tipo de interés común. En este contexto, estaremos siempre en presencia de racionamiento de crédito, es decir, las curvas de oferta de crédito y de demanda de crédito nunca serán tangentes por lo que el equilibrio siempre se dará en un lugar sub-óptimo determinado por la función de beneficio-riesgo del banco (en función de su calidad de cartera). Eésta función se puede expresar como:

\eqblock{
\begin{aligned}
    &\max_{\{i\}}{r^e(i)=P(\cdot |i)(1+i)+(1-P(\cdot|i))c}\\[2em]
    \text{Donde,}
    &\begin{cases}
        r^e&\equiv \quad \text{Rendimiento esperado por el banco en función de $i$}\\[0.5em]
        P(\cdot|i)&\equiv \quad \text{Probabilidad de repago condicionada a $i$}\\[0.5em]
        (1+i)&\equiv \quad \text{Rendimiento bruto: capital más intereses}\\[0.5em]
        c&\equiv \quad \text{Recuperación esperada en caso de impago (garantía)}\\
    \end{cases}
\end{aligned}
}{Fricciones reales: Mercado de crédito. STIGLITZ-WEISS}

Como consecuencia de esta formulación (derivando y solucionando para la CPO) obtendremos que el tipo de interés fijado por el banco no equilibra la oferta de crédito con la demanda de crédito ya que el rendimiento esperado ($r^e$) no es una función monótamente creciente en $i$. De hecho, será aquel tipo de interés que no expulse del mercado a los clientes de bajo riesgo, de tal formaque se puede obtener una cartera de crédito relativamente estable (puesto que también se incluirán los prestatarios de alto riesgo a un precio más moderado, ya que el tipo de interés será inferior al óptimo en su caso). Gráficamente este tipo de interés será:

\imagenfit{Mercadodecredito_Pooling.png}{Reducción del crédito. Modelo STIGLITZ-WEISS}{Apuntes ICEX-CECO. Tema 3.A.7}

Si ilustramos gráficamente como serían la curva de oferta de crédito ($S_c$) y la curva de demanda de crédito ($S_d$) como función del tipo de interés, tenemos:

\imagenfit{MercadoCredito_Racionamiento.png}{Racionamiento del crédito y shocks de Demanda Agregada}{Apuntes Víctor Begué}

Además, ante un shock de demanda agregada, el efecto de una política monetaria expansiva desplazaría la curva de oferta de crédito verticalmente (hacia arriba), reduciendo el gap debido al racionamiento, aumentado, para un mismo nivel de tipos, la oferta total de crédito y el rendimiento esperado del prestatario. Tratándose por tanto de un ajuste vía cantidad y no vía precios.

\paragraph{Mercado de trabajo}
Por último, podemos también estudiar como interactúan las rigideces nominales y reales en el mercado de trabajo, conduciendo a la existencia, en equilibrio, de desempleo involuntario, uno de los pilares fundamentales de la primera generación de econosmitas de la NEK.

\textsc{Modelo de SOLOW: Salarios de eficiencia}

Uno de los ejemplos más sencillos para ilustrar esta interacción es el Modelo de Salarios de eficiencia de SOLOW (1979):\footnote{Estas teorias surgen en los años 80's y se centran en que los salarios no se determinan solo por la interacción de oferta y demanda sino que las empresas tienen incentivos a pagar salarios superiores a los de vaciado de mercado para aumentar la productividad y la eficiencia de los trabajadores. 
} la existencia de información asimetrica post-contractual relacionada con el fenómeno de riesgo moral puede explicar la fijación de salarios superiores a los óptimos (imperfecciones reales conducen a fricciones nominales). El principal motivo seria que el esfuerzo del trabajador no es observable por lo que habrá que incentivar su esfuerzo con un salario más elevado. La empresa se enfrenta al problema de fijar el empleo y el salario que va a pagar a los trabajadores de forma que maximice su beneficio. El output dependerá entonces tanto del número de trabajadores como de su esfuerzo, y este esfuerzo depende, a su vez, del salario. Por tanto, para aumentar el esfuerzo y el output, se fijará un salario mayor que el de vaciado de mercado.

La condición de optimalidad que determina el nivel salarial (nominal) que induzca en el trabajador un nivel de esfuerzo que optimice los beneficios de la empresa se obtiene:

\eqblock{
\begin{aligned}
    &\max_{\pi_i}{\pi_i=F(e(w)\;\cdot\;L)-w\;L}\\[1em]
    &\text{De las CPO's:}\quad \frac{\partial \pi_i}{\partial L}=\frac{\partial \pi_i}{\partial w}\Longrightarrow\underbrace{\boxed{\frac{w\;\cdot\;e'(w)}{e(w)}=1}}_{\substack{\text{\scriptsize Elasticidad del esfuerzo respecto al }\\\text{\scriptsize salario es igual a $1$}}}
\end{aligned}
}{Fricciones reales-nominales: Mercado de trabajo. Modelo de SOLOW}

Por tanto, como vemos, con el salario que resuelve la ecuación se minimiza el coste del trabajo efectivo (definido como esfuerzo por trabajo): cuando la empresa contrata a un trabajador, obtiene $e(w^*)$ unidades de trabajo efectivo a un precio $w^*$, con lo que el coste unitairio del trabajo efectivo será $w^*/e(w^*)$, y cuando la ecuación se iguale ($w^*$), el coste del trabajo efectivo es mínimo y la empresa maximiza sus beneficios. Gráficamente:

\imagenfit{MercadoTrabajo_SOLOW.png}{Salarios de eficiencia: Equilibrio}{Apuntes ICEX-CECO. Tema 3.A.7}

\subsection{Recepción} 
Tal y como se ha señalado con anterioridad, las aportaciones la NEK son keynesianas en cuanto que justifican la existencia de rigideces de precio y el desempleo involuntario pero no realiza una defensa incondicional de las politica de demanda activas. Entre los autores de la NEK existen opiniones muy diversas pero hay algunas conclusiones comunes. 

\subsubsection{Aportaciones}
\paragraph{No-neutralidad del dinero}
Los autores de la NEK le prestan una mayor atención a la politica monetaria en tanto en cuanto la NEK surge como contraposición a la inefectividad de las politicas de estabilizacion que propugnaba la NMC. 

En general, los autores de la NEK coinciden en que la politica monetaria tiene efectos reales a corto plazo por la inercia del nivel general de precios que hace que los shocks nominales influyan sobre la producción y el empleo. 

Pero, a largo plazo, los efectos se disipan con el tiempo con lo que la economía vuelve al nivel de produccion previo pero con un mayor nivel de precios. En este sentido, introducen el modelo agregado de la 1\textsuperscript{era} NEK: NAIRU. Por tanto, en general recomiendan usar la politica monetaria solo para controlar la inflacion porque al tener efectos reales un uso continuado de la misma genera inestabilidad.\footnote{Segun BALL, MANKIW y ROMER (1988) si la politica monetaria se usa con fines expansivos, la tasa media de inflación aumenta y disminuye la eficacia de la politica monetaria y aumentan los precios.} 

\paragraph{Ciclos reales. Política económica: discrecionalidad}
En relacion con los ciclos, los autores de la NEK se basan en la información imperfecta, las expectativas racionales y la competencia imperfecta junto con rigideces de precios lo que provoca que se genere un ajuste de las cantidades: ante una perturbación de la Demanda Agregada, ya sea fiscal o monetaria, el lento ajuste de los precios en un contexto de competencia imperfecta supone un mecanismo de transmisión de este shock al output y el empleo. Por tanto, los shocks de la demanda agregada tienen efectos reales. 

Para éstos autores, las causas más habituales de los shocks son las pertubaciones producidas por las politicas fiscales o monetarias.\footnote{Por ejemplo, los cambios anticipados o no en la politica monetaria pueden causar fluctuaciones en la producción. Esto es una distinción importante respecto de los modelos neoclasicos de los ciclos basados en las expectativas racionales en los que los mercados se ajustan instantáneamente y solo los cambios no anticipados de política monetaria llevan a fluctuaciones de la producción (LUCAS, 1972).
} Sin embargo, también pueden provocar fluctuaciones las perturbaciones exógenas del consumo y la inversión derivadas de cambios en el grado de optimismo o pesimismo en la economía, especialmente cuando dichos sentimientos no están avalados por las condiciones económicas fundamentales.\footnote{Además, estos sentimientos o expectativas pueden llevar a situaciones de profecías o expectativas que se autocumplen (Woodford 199 1 y Shleitler 1986). Es decir, dichos modelos explican tanto los ciclos inducidos por cambios en la política económica como por cambios en las expectativas o en el sentimiento de los agentes económicos. Dichas fluctuaciones económicas llegan a su fin, bien porque la política económica sabe rectificar y responder adecuadamente a dichos desequilibrios o porque el sentimiento o las expectativas de los consumidores o inversores cambia y se ajusta, independientemente de la política económica, en respuesta a los desequilibrios.
}

A pesar de esto, respecto al debate sobre las reglas y la discrecionalidad, los autores de la NEK aceptan y defienden la discrecionalidad. 

\subsubsection{Limitaciones}
Los trabajos de esta primera oleada de economistas de la NEK no fueron capaces de frenar el desarrollo de la NMC como rama dominante en la macroeconomía moderna. Según DE VROEY (2016), este fracaso puede deberse a dos motivos principales:
\begin{lnum}
    \item Por un lado, tuvieron una aproximación defensiva, anclados en metodologías y supuestos que se estaban volviéndose obsoletos (por ejemplo, la consideración de contextos estaticos y de equilibrio parcial); y,
    \item Por otro lado, aunque competían con las aportaciones de LUCAS en el marco conceptual y teórico, los desarrollos aplicados de la TCR los relegaron a una posición únicamente académica y fuera del ámbito de la guía de política económica.
\end{lnum} 

\clearpage
\section{Segunda generación}
Una década despues, la llegada de la segunda oleada de economistas neo-Keynesianos cambió de nuevo el terreno de juego y supuso una nueva evolución en la disciplina macroeconómica. Éstos autores aceptan las reglas de juego propuestas por la NMC pero cambiarán la esencia de los modelos de la Teoría del Ciclo Real. 

\imagenfit{Dif_NEK1yNEK2.png}{Diferencias entre las dos generaciones}{Apuntes ICEX-CECO. Tema 3.A.7}

\subsection{Contexto histórico}

\subsubsection{Autores y Universidad de referencia}

\subsection{Método}
Básicamente, se aceptan los principios definitorios de los modelos DSGE: agentes racionales, cuyas decisiones estan microfundamentadas en un contexto dinamico y de equilibrio general. Sin embargo, se introducen dos cambios muy relevantes: 
\begin{lnum}
    \item Por un lado, se rechaza la existencia de competencia perfecta y se introducen marcos de competencia monopolística con rigideces y ajustes escalonados de precios y salarios; y, 
    \item Por otro, se postula la vuelta a primera plana de los aspectos monetarios y de la politica monetaria.
\end{lnum} 

Estos elementos novedosos introducidos por la NEK en el modelo TCR son de gran relevancia. Por un lado, la NEK de segunda generación incorpora la competencia monopolística como marco de análisis.\footnote{Si bien el concepto fue introducido por CHAMBERLAIN (1933), no tuvo una mayor difusión hasta que fue revisado por DIXIT y STIGLITZ en 1977. Desde entonces fue ampliamente usado en modelos de comercio y crecimiento. Los neokeynesianos lo incorporaron de lleno al campo de la macroeconomía.
} Por otro lado, trajeron de nuevo las cuestiones monetarias y el papel del dinero al centro del debate económico. 

\subsubsection{Apalancamiento sobre la Nueva Macroeconomía Clásica}
Por tanto, el modelo base de la NEK es un desarrollo del modelo basico de la TCR que, como este: parte de un agente representativo que resuelve un problema de optimización dinámica dados unos recursos limitados, quedando esta decisión caracterizada por la ecuación de Euler.\footnote{La ecuación de Euler establece que la relación marginal de sustitución entre consumo presente y consumo futuro ha de igualarse al precio relativo del consumo presente (dado por el tipo de interes). 
} Sin embargo, el contexto real es ahora el de competencia monopolística con rigideces nominales temporales de precios y una autoridad monetaria, de modo que el equilibrio es mas complejo al tener que aunar:
\begin{la}
    \item Un agente representativo; 
    \item Las empresas (producción); y,
    \item La propia autoridad monetaria.
\end{la} 

Por todo esto, tal y como señalan GOODFRIEND y KING, sería más propio decir que esta nueva ola de modelos DSGE constituye una Nueva Sintesis, comunmente aceptada hasta la llegada de la Gran Recesión. 

Las principales caracteristicas de los modelos de equilibrio general dinamico de la NEK podrian resumirse en los siguientes puntos (Gali, 2002 y Gali, 2007):
\begin{lnum}
    \item \textit{Modelos DGSE}. Instrumentos derivados de la TCR, en particular, los modelos DGSE que se basan en el comportamiento optimizador por parte de los hogares y las empresas y el uso de expectativas racionales y vaciado de mercados;
    \item \textit{Competencia monopolítica}. Las empresas se modelizan en un marco de competencia monopolistica;
    \item Rigideces nominales, que son un elemento central del modelo y una fuente de la no neutralidad monetaria. Generalmente se introducen en la forma de restricciones en la frecuencia de fijacion de los precios y salarios. La rigidez de precios es la fuente de no-neutralidad del dinero en el corto plazo;
    \item Las decisiones de fijacion de precios se vuelven forward-looking, ya que los individuos reconocen que los precios fijados hoy permanecen mas allá de un periodo. Por tanto, existe una nueva perspectiva de la naturaleza dinámica de la inflación de tipo forward-looking. La decision de las empresas sobre su precio depende de los costes esperados y demanda futuras y se refleja en la Curva de Phillips de la NEK;
    \item \textit{Output gap}. Este concepto es central en los modelos NEK. Se define como la desviación del ouput con respecto a su nivel de equilibrio en ausencia de rigideces nominales;
    \item \textit{Mecanismo de transmision}. El tipo de interés es el mecanismo de transmision de los shocks de Polltica Monetaria a las variables reales y no necesariamente a través del efecto liquidez keynesiano;
    \item \textit{Bienestar basado en la utilidad.} El modelo de la NEK permite hacer un análisis de bienestar basado en la utilidad de las consecuencias de las políticas monetarias alternativas; ademas, permite diseñar una política monetaria óptima. La política monetaria óptima permite estabilizar el nivel de precios y el output gap completamente;
    \item \textit{Regla de politica monetaria óptima}. Una buena aproximación de la regla de la política monetaria óptima es una regla de tipos de interés (regla de política simple). Esta regla consiste en el ajuste del tipo de interés por parte del Banco Central en respuesta al aumento de la inflación o del output gap. En este escenario la cantidad de dinero, $M$, es endógena; y,
    \item Aporta argumentos respecto al debate reglas vs discrecionalidad y en politica monetaria óptima investigan sobre si existe \textit{trade-off} entre el output y la inflación.
\end{lnum}

\subsection{Objeto: Modelo de tres ecuaciones}
Los nuevos modelos de la NEK-2  sintetizan las condiciones de optimización mediante tres elementos:\footnote{Las dos primeras ecuaciones son una log-linearización alrededor del Estado Estacionario de las reglas de decisión del agente representativo y de las empresas respectivamente. La tercera ecuación describe la regla de política monetaria óptima del Banco central.
} 
\begin{lnum}
    \item La demanda agregada (curva IS ampliada por las expectativas);
    \item La curva de Phillips ampliada por las expectativas; y,
    \item La regla de Taylor. 
\end{lnum}

El resultado es un sistema de ecuaciones que podria considerarse una modelización del modelo IS-LM. Este tipo de modelos sirven para explicar los ciclos económicos y para analizar la relación entre dinero, inflación, ciclo y rigideces así como para evaluar las distintas opciones de política económica.

\subsubsection{Curva IS ampliada por expectativas}
La primera ecuacion del modelo seria la curva IS am pliada por las expectativas, que proviene de la derivaci6n del problema de maximizacion intertemporal de la utilidad, que depende del consumo y el ocio de los hogares, sujetos a una restriccion presupuestaria. Los bienes de consumo estan diferenciados y los hogares consumen distintas variedades de los mismos. Los hogares toman decisiones de consumo, oferta de trabajo y de inversion en titulos de deuda publica (bonos).

El equilibrio, queda caracterizado por las condiciones habituales:
\begin{lnum}
    \item \textit{Condición de sustitución intertemporal del consumo} (ecuación de Euler). Igualdad entre la $RMS$ de consumo entre períodos y el precio relativo de consumo entre períodos (uno más la tasa de inflación);
    \item \textit{Condición de sustitución intertemporal del ocio}. Igualdad entre la $RMS$ de ocio entre períodos y el salario nominal relativo entre períodos (uno más inflación salarial); y,
    \item \textit{Condición de sustitución intratemporal de consumo-ocio}. Igualdad entre la RMS consumo-ocio y el salario real.
\end{lnum}

La curva (log-linearizada) IS queda entonces definida como:

\eqblock{
\begin{aligned}
    x_t=-\frac{1}{\theta}\;\left(\overbrace{i_t-E_t[\pi_{t+1}\;|\;\Omega_t]}^{\substack{\text{\scriptsize Tipo de interés real: $r_t$}\\ \text{\scriptsize (ecuación de FISCHER)}}}-\underbrace{r_t^n}_{\substack{\text{\scriptsize Tipo de interés}\\\text{\scriptsize natural}}}\right)+E_t[x_{t+1}\;|\;\Omega_t]+u_t
\end{aligned}
}{Curva IS dinámica o ampliada con expectativas. Modelo NEK-2}

La curva IS ampliada por las expectativas representa el lado de la demanda de la economía.\footnote{Esta curva depende de:
\begin{lnum}
    \item $r_n^t$. El tipo de interés natural es el valor que tomaría el tipo de interés real consistente con el output en su nivel de equilibrio con precios flexibles (en ausencia de rigideces), si no existieran distorsiones. Es decir, es el tipo de interés de equilibrio competitivo que sería eficiente y en el que se cumple la hipotésis de neutralidad del dinero. En otras palabras, el tipo de interés natural es independiente de la política monetaria y solo depende de variables estructurales de la economía:
    \begin{la}
        \item \textit{Positivamente} del tipo de descuento intertemporal, $\rho$, pues cuando los agentes se vuelven más impacientes tienden a ahorrar menos, aumentando el coste de la financiación; y,
        \item \textit{Negativamente} de los shocks tecnológicos (y positivamente de shocks de demanda, $u_t$) ya que, un aumento en el output potencial aumenta tanto el consumo como el ahorro, empujando hacia abajo el tipo de interés natural.
    \end{la}
    \item Tipo de interés real (rt), compuesto por: (i) tipo de interés nominal ($i_t$); y, (ii) las expectativas de inflación futura ($E_t[\pi_{t+1}|\;\Omega_t]$);
    \item Elasticidad de Sustitución Intertemporal ($1/\theta$).
    \item Output gap ($x_t$), que es la desviación del output observado de la tasa natural: el que existiría si hubiera perfecta flexibilidad de precios y salarios y el que se buscaría recuperar en caso de perturbaciones por la autoridad monetaria;
    \item Expectativas de output gap futuro ($E_t[x_{t+1}\;|\;\Omega_t]$); y,
    \item $u_t$ un shock de demanda, $u_t\sim N(0,\sigma_x^2)$.
\end{lnum}
Véase: \href{https://research.stlouisfed.org/publications/economic-synopses/2018/09/28/measuring-potential-output}{Measuring output gap}, para más información sobre la Curva IS dinámica.
} A diferencia de la curva IS tradicional de la sintesis neoclásica, se trata de una relación que incorpora las expectativas, compatible con la HER y, establece que, por el lado de la demanda existe una relación de equilibrio entre el output gap y el tipo de interés nominal, condicionada a las expectativas sobre la evolución futura de precios y output gap:
\begin{lnum}
    \item Si aumenta el output gap esperado en el futuro, aumentará el consumo y la demanda agregada hoy, por lo que aumentará el output gap de hoy para cualquier nivel de tipos de interés;
    \item El término $-1/\theta$ mide el impacto del aumento de los tipos de interés reales sobre la demanda agregada; y,
    \item Por su parte, la diferencia entre los tipos de interés reales y los naturales ($i_t-E_t[\pi_{t+1}\;|\;\Omega_t]-r_t^n$), introduce un mecanismo a través del cual se transmiten las perturbaciones monetarias: si el tipo de interés nominal aumenta, en presencia de rigideces de precios, el tipo de interés real se elevará por encima del natural, lo que llevará a los hogares a posponer su consumo, reduciendo el consumo presente e incrementando el consumo futuro.
\end{lnum}

Así, se pueden extraer dos conclusiones de la curva: 
\begin{lnum}
    \item Por un lado, como $\theta$ es positivo, cuando se espera una bajada de los precios, los agentes posponen sus decisiones de consumo al siguiente período y el output gap baja; y,
    \item Por otro lado, el banco central puede afectar a las variables reales modificando los tipos de interés nominales a corto plazo. Esto provoca un cambio en los tipos de interés reales lo que lleva a los agentes a cambiar sus decisiones de optimización.
\end{lnum} 

\subsubsection{Curva de PHILLIPS de la NEK}
El segundo elemento del modelo canónico es la curva de Phillips de la NEK, donde subyacen elementos importantes de la NEK: competencia monopolística y rigidez de precios (sticky prices), junto con las hipótesis de comportamiento optimizador de los agentes y expectativas racionales.\footnote{En la primera mitad de los años 70, los autores de la NMC, a través de la función de oferta agregada de LUCAS (1972) criticaron las políticas activas de estabilización keynesianas ya que, con base al supuesto de equilibrio continuo de los mercados y la HER, llegaron a la conclusión de la inefectividad de las políticas de estabilización sistemáticas, que no tienen incidencia alguna sobre el nivel de producción real o sobre la tasa “natural” de desempleo inclusive en el corto plazo. Además, también concluyeron que los componentes aleatorios no sistemáticos de las políticas fiscales y monetarias solo aumentan la incertidumbre y las fluctuaciones cíclicas. El elemento clave de la NMC fue la hipótesis de equilibrio continuo de los mercados.

La primera reacción de los economistas de inspiración keynesiana consistió, tal y como hemos visto, en levantar la hipótesis de flexibilidad perfecta de precios y salarios. Entre los primeros estudios al respecto destacan los de STANLEY FISCHER (1977) y JOHN TAYLOR (1980), quienes estudiaron cómo afectaba el escalonamiento de precios y salarios en la economía. Estos estudios fueron complementados por los desarrollos de JULIO J. ROTEMBERG (1982) y GUILLERMO A. CALVO (1983).
}

Se suele usar para definir la Curva de PHILLIPS de la NEK el modelo de fijación de precios à la CALVO de rigideces nominales de precios en el bien final. En cada período, solo hay un determinado porcentaje de empresas, $1-\theta$, que pueden ajustar sus precios, mientras que el resto de empresas, $\theta$, no pueden ajustarlos y tienen que mantener los del período anterior.\footnote{De este modo, el grado de rigidez de precios de la economía se puede medir por el valor de $\theta$: cuanto mayor sea esta porción, menos empresas pueden ajustar los precios y, por lo tanto, mayor será el grado de rigidez de precios, ya que el tiempo que transcurrirá entre cambios de precios será más elevado. 
} Las empresas que pueden variar sus precios los determinarán óptimamente, maximizando el valor descontado de los beneficios. Además, se tiene en cuenta que, cuando las empresas ajustan sus precios, tienen en cuenta la existencia de la rigidez nominal en su decisión, ya que es posible que en el futuro deseen ajustar los precios y no puedan debido a ésta. 

De acuerdo con la formulación de CLARIDA, GALÍ y GERTLER (1999), la solución óptima para la empresa sería fijar el precio igual a una media ponderada de los precios que se esperarían en el futuro si no hubiera fricciones, por tanto, bajo el supuesto de que existe competencia monopolística, este sería fijar un mark-up ($\mu$) sobre el coste marginal ($CMg$):\footnote{
De acuerdo con la ecuación, la tasa de inflación observada dependerá de:
\begin{lnum}
    \item La inflación esperada para el siguiente periodo;
    \item El gap entre precio óptimo sin fricciones y el nivel deneral de precios actual; y,
    \item Una variable estocástica de shocks en costes $e_t\sim N(0, \sigma_c^2)$
\end{lnum}
}

\eqblock{
\begin{aligned}
    &P^*=CMg_t+\mu\\[0.5em]
    &\pi_t=\beta\;E_t[\pi_{t+1}\;|\;\Omega_t]+\left(\frac{(1-\theta)\;(1-\beta\;\theta)}{\theta}\right)\;(\mu+CMg_t-P_t)+e_t\\[1em]
    \text{Reformulando:}&\\[0.5em] 
    &\mu+CMg_t-P_t\to \underbrace{dCMg=\mu+CMg_t-P_t}_{x_t}\quad \text{y,}
    \begin{cases}
        \text{Utilizando output gap como proxy}\\[0.5em]
        \left(\frac{(1-\theta)\;(1-\beta\;\theta)}{\theta}\right)=\delta
    \end{cases}\\[1em]
    &\Longrightarrow \underbrace{\boxed{\pi_t=\beta E_t[\pi_{t+1}\;|\;\Omega_t]\;\delta x_t+e_t}}_{\text{Curva de PHILLIPS de la NEK}}
\end{aligned}
}{Obtención de la Curva de PHILLIPS de la NEK. Modelo NEK-2}

Un problema al tratar de implementar el modelo empíricamente es que los datos de los costes marginales reales no son observables. Los sistemas de cuentas nacionales contienen información en relación a los factores que afectan a los costes marginales como los salarios, pero no del coste concreto de producir una unidad más de output. Además, es muy
probable que los costes marginales sean procíclicos y mucho más que los precios. De este modo, cuando los niveles de producción son relativamente altos en relación al output potencial, hay más competencia para conseguir los factores de producción disponibles, y esto tiende a incrementar los costes. Por ejemplo, aumentan los costes de los factores por encima de los incrementos en los precios. Por estos motivos, la curva de Phillips se suele reformular usando el output gap ($x_t$) como un proxy de la desviación del coste marginal.

\paragraph{Críticas a la Curva de PHILLIPS de la NEK}
La curva de Phillips de la NEK ha sido criticada por diversos motivos, ya que predice una dinámica difícilmente compatible con la realidad. Entre las más destacadas:
\begin{lnum}
    \item Según BALL (1994), en los modelos que usan reglas de fijación de precios à la CALVO, el anuncio de una política de desinflación creíble que se va a implementar en el futuro provocaría, de forma paradójica, una expansión en lugar de una recesión, ya que la inflación esperada disminuirá y dado que la inflación actual se forma con los precios esperados se reducirá la inflación actual de forma inmediata. El efecto inmediato será un aumento de los saldos reales que tendría un efecto expansivo. Sin embargo, la evidencia empírica muestra que la desinflación ocurre gradualmente y no de forma repentina y, además, se asocia con un crecimiento moderado y un desempleo alto; y,
    \item Además, la curva de PHILLIPS de la NEK es incapaz de justificar la persistencia de la inflación y el efecto gradual y retardado de la política monetaria sobre la inflación. Por otro lado, se ha apuntado a que generalmente la inflación pasada es un mejor predictor de la inflación futura que la inflación esperada. Ante estas críticas, se han realizado reformulaciones de la curva de PHILLIPS para introducir una cierta inercia en la inflación.
\end{lnum}

\paragraph{Modificaciones: Curvas de PHILLIPS \textit{híbridas}}
Entre otras propuestas, también cabe destacar modelos que consideran una indicación generalizada (CHRISTIANO et al., 2005) o reglas de ajuste sencillas (GALÍ y GERTLER, 1999) y que conducen a curvas de PHILLIPS del tipo:

\eqblock{
\begin{aligned}
    \pi_t=\alpha\;\pi_{t-1}+\beta\;E[\pi_t+1\;|\;\Omega_t]+\delta\;x_t+e_t
\end{aligned}
}{Curva de PHILLIPS híbrida. Modelo NEK-2}

Estos autores parten del modelo de CALVO en el que una proporción de las empresas podían cambiar los precios de manera óptima en cada período y asumen que el resto de empresas en lugar de mantener sus precios inalterados, los actualizan conforme a la inflación pasada. Este hecho provoca que la inflación presente también dependa tanto de la inflación desfasada como de la inflación esperada en el futuro, además de ser función del output gap.\footnote{Este tipo de modelos superan algunos de los problemas asociados con la curva de Phillips de la NEK pero siguen habiendo cuestiones pendientes en relación a su capacidad para explicar los cambios en la inflación y, sobre todo, fallan cuando se contrastan empíricamente, ya que llegan a la conclusión de que todas las empresas cambian sus precios aunque no sea óptimo. Es más, con indexación, y pocas empresas fijando sus precios de manera óptima, estos modelos sugieren que, si bien los Bancos Centrales pueden provocar un cambio en los tipos de interés reales de manera sencilla, se requieren grandes movimientos en los tipos de interés para estabilizar la inflación, generando grandes oscilaciones de los tipos a lo largo del ciclo.
}

Por su parte, MANKIW y REIS (2002) plantean sustituir la rigidez o inercia de los precios por la inercia de la información, ya que su difusión es más lenta.\footnote{El supuesto que hacen es que existen costes de información y las empresas racionan la información que adquieren para formar sus expectativas. En cada período, todos los agentes actualizan sus precios, pero el problema es que necesitan realizar una previsión de la inflación para hacer esto. Antes de realizar su previsión, las empresas necesitan actualizar su información para incluir los últimos datos disponibles y el modelo supone que solo existe una fracción de agentes que actualiza la información disponible necesaria para actualizar esos precios, el resto no puede actualizarla. Es decir, solo toman una parte de sus decisiones basándose en la información disponible. Por tanto, los precios sí pueden cambiar rápidamente pero solo una fracción de los mismos se basa en información actualizada. El porcentaje de empresas que actualizan la información es aleatorio por lo que los shocks pasan por los precios de manera gradual porque las empresas tardan un tiempo en darse cuenta de que ha ocurrido un shock. Según este modelo, las expectativas formadas en el pasado sobre los precios actuales tienen un papel en la dinámica de la inflación.
}

Pese a las distintas críticas, la curva de Phillips de la NEK sigue siendo útil para cuestiones de política económica y para cuestiones pedagógicas, pues consigue explicar la persistencia de la inflación y porque responde de manera gradual a los shocks.

\subsubsection{Regla de TAYLOR: Autoridad monetaria}
Como se ha visto, la curva de Phillips neokeynesiana determina la inflacion dada una senda de output-gap. Por su parte, a partir de IS dinamica, el output gap queda determinado por la senda de tipos de interes. Por lo tanto, para cerrar el modelo es necesaria una ecuación que describa la evolucion del tipo de interés nominal, es decir, una ecuación que describa el comportamiento de la politica monetaria. 

\imagenfit{CirculoNEK_2.png}{La regla de política monetaria cierra el círculo del Modelo de la NEK-2 de tres ecuaciones}{Apuntes Víctor Gutiérrez. Tema 3.A.7}

Los modelos de la NEK se basan en la instrumentación de la política monetaria a través de una regla de política monetaria. Dentro de éstas, la que más éxito ha tenido entre la NEK es la regla de TAYLOR.\footnote{La regla surge de una propuesta realizada por TAYLOR en una conferencia en 1992 (TAYLOR, 1993). TAYLOR introdujo la regla como una regularidad empírica más que como una conjetura teórica e incidió en que la regla se ajustaba a la política monetaria que se habia aplicado en los últimos años. La regla no tardó en generalizarse y aplicarse como una prescripción de política económica. La regla propuesta por TAYLOR originalmente fue: $i_t=2\%+\pi_t+0,5(y_t−y^*)+0,5(\pi_t−2\%)$.

Es decir, $\phi_y=\phi_{\pi}=0,5$ y el objetivo de inflación y los tipos reales eran el 2\%
} Esta se basa en el tipo de interés nominal a corto plazo del banco central ($i_t$), que es usado como instrumento de política económica. De acuerdo con la regla, debería ser fijado como respuesta al output gap o diferencia entre el PIB real y el potencial, $x_t=(y_t-y^*)$, y al gap de la inflación, entendido como la diferencia entre la tasa de inflación observada y la deseada ($\pi=\pi_t-\pi^*$).

\eqblock{
\begin{aligned}
    i_t=r_t^*+\pi_t\;\phi_y(y_t-y^*)+\phi_\pi(\pi_t-\pi^*)
\end{aligned}
}{Regla de TAYLOR. Modelo NEK-2}

De este modo, la regla indica que si no hay ninguna desviación del output gap, ni de la inflación respecto a su objetivo, el tipo de interés nominal se igualará al tipo de interés real de equilibrio ($r_t^*$), más la inflación. Si se producen desviaciones positivas en el output gap o en la inflación respecto al objetivo, el banco central debería incrementar los tipos nominales.\footnote{La regla presenta numerosas ventajas, la principal es que es simple. La segunda es que es util para los Bancos centrales a la hora de tomar sus decisiones. Ademas, su ajuste empirico es excelente. No obstante, tambien presentan algunas desventajas, como su caracter backward-looking, el hecho de que no está microfundamentada y que no tiene un componente no estocástico.
} 

En el modelo canónico, se puede usar una regla tipo TAYLOR modificada de la siguiente forma: 

\eqblock{
\begin{aligned}
    i_t=v\;i_{t-1}+\phi_\pi\pi_t+x_t+\varepsilon_t
\end{aligned}
}{Regla de TAYLOR (suavizada). Modelo NEK-2}

Donde $v$ indica el grado de suavización del tipo de interes, $\phi_\pi$ es el peso relativo que le asigna el Banco central al gap de inflación y $\varepsilon_t$ es un shock monetario (que puede ser definido como cambios en la política monetaria que no son el resultado de respuestas lógicas de la autoridad monetaria en respuesta a movimientos de otras variables). La autoridad monetaria tiene que ajustar el tipo de interés más que proporcionalmente respecto al cambio en la inflación para asegurar la senda adecuada de los tipos reales. De modo que si el Banco central actúa de esta manera, el modelo genera un equilibrio estable y único.

\subsection{Aportaciones}
La importancia de las aportaciones de la NEK-2 y la coincidencia en el instrumental utilizado con respecto a la TCR ha hecho que, como hemos comentado antes, sea frecuente hablar de una \textit{Nueva Síntesis} entre ambas escuelas ya que, tras años de duro debate en el seno de la macroeconomía, se volvieron a encontrar puntos en común sobre algunas cuestiones metodológicas (como los modelos DSGE) y sobre cuestiones de política económica.

\imagenfit{NuevaSintesis.png}{Elementos característicos y descartados de la \textit{Nueva Síntesis}}{Apuntes ICEX-CECO. Tema 3.A.7}

A partir del modelo canónico, se pueden extraer distintas conclusiones de política económica que han fueron comúnmente aceptadas durante los años 90 y los 2000, hasta el inicio de la crisis de 2008. 

\subsubsection{Política monestaria}
Si bien se plantea que a largo plazo la política monetaria no tiene efectos reales, a corto plazo, las rigideces nominales crean una relación cíclica entre output e inflación. Además, los resultados de la política mejoran con un banco central independiente, pues el manejo de las expectativas es esencial. Sin embargo, como veremos, para los autores de la NEK-2 la inflación tiene numerosos costes y, dada la curva de PHILLIPS, su control y estabilidad es esencial para conseguir que el output crezca de acuerdo con su nivel potencial en el largo plazo.

\paragraph{Deseabilidad del Equilibrio: Principio de TAYLOR}
A la hora de evaluar el desempeño de una regla de política monetaria, la estabilidad dinámica que ésta pueda aportar a la economía juega un papel fundamental, siendo deseable disponer de una regla de actuación del banco central que genere un único equilibrio, estable y que nos permita alcanzar el máximo nivel de bienestar posible. Es por ello que se realizan análisis de estabilidad sobre distintas reglas con el fin de estudiar el comportamiento global o local de los modelos estructurales bajo las mismas. De acuerdo con HEIJDRA y VAN DER PLOEG (2002), un equilibrio estable se define como aquél en el que el equilibrio único es restaurado eventualmente tras la materialización de un shock durante uno o más períodos (también conocido como estado estacionario).

De hecho, como los modelos pueden ser susceptibles tanto de tener equilibrios múltiples como equilibrios inestables o divergentes, BLANCHARD y FISCHER (1993) ofrecen una clasificación de los tipos de equilibrios múltiples a evitar en el diseño de éstas:
\begin{lnum}
    \item \textit{Equilibrios burbuja}. En estos modelos existen componentes cuyo valor explota en términos esperados a lo largo del tiempo. MCCALLUM (2004) enfatiza que una regla de política económica bien diseñada debe conducir a un único equilibrio caracterizado por la ausencia de burbujas; y,
    \item \textit{Equilibrios de manchas solares}. En estos modelos surgen equilibrios múltiples donde las variables endógenas cambian a estados aleatorios, no asociados a ningún cambio en el valor de los parámetros de fondo o estructurales. Por ejemplo, los equilibrios de profecías autocumplidas.\footnote{Una profecía autocumplida, autorrealizada o, mejor, autorrealizable es una predicción que, una vez hecha, es en sí misma la causa de que se haga realidad.\\
    En el siglo XX la expresión fue acuñada por el sociólogo R. K. MERTON, quien formalizó su estructura y sus consecuencias, en su libro \textit{Teoría social y estructura social}, da la siguiente definición: \textit{La profecía que se autorrealiza es, al principio, una definición «falsa» de la situación, que despierta un nuevo comportamiento que hace que la falsa concepción original de la situación se vuelva «verdadera»}.
    }
\end{lnum}

En el modelo de la NEK-2 planteado, la condición matemática necesaria para garantizar que existe un único equilibrio es la siguientes: 

\eqblock{
\begin{aligned}
    \text{Condición de unicidad:}\qquad \underbrace{\boxed{\phi+\frac{1-\beta}{\delta}>1}}_{\text{Garantiza el Principio de TAYLOR}}
\end{aligned}
}{Condición de unicidad del Equilibrio. Modelo NEK-2}

Esta condición garantiza que se cumpla el Principio de TAYLOR, es decir, que un incremento de inflación sea respondido con un incremento del tipo de interés nominal proporcionalmente mayor. 

El Principio de Taylor suele ser condición suficiente para garantizar la unicidad del equilibrio del modelo. La intuición de la condición obtenida es la siguiente: un incremento en la expectativa de inflación genera a su vez un incremento del output gap de largo plazo de $(1-\beta)/\delta$, lo que incrementa a su vez las presiones inflacionistas presentes. Por tanto, los coeficientes de la Regla de TAYLOR deben de cumplir conjuntamente la condición para garantizar que el tipo de interés nominal responde de forma más que proporcional a incrementos en el nivel de inflación.\footnote{
Por ejemplo, los coeficientes estimados por Taylor ($\phi_\pi=1.5$ y $\phi_y=0.5$) satisfarían este criterio.

Other insights generated by the New Keynesian model pertain to its normative implications for the conduct of monetary policy. One finding in that regard has to do with the insight that if the central bank applies a rule which adjusts the policy interest rate sufficiently strongly in response to variations in inflation and output (a condition known as the Taylor principle) then the economy will have a unique equilibrium. Otherwise, the equilibrium is locally indeterminate, opening the door to fluctuations driven by self-fulfilling revisions in expectations (sometimes known as “sunspot fluctuations”). Clarida et al. (2000) provide evidence suggesting that the local uniqueness condition may not have been satisfied during the pre-Volcker era, potentially giving rise to unnecessary instability and providing an explanation for the macroeconomic turbulence of that period.
}

\paragraph{Deseabilidad del Equilibrio}
Una vez hemos demostrado que la Regla de Taylor nos conduce a un equilibrio único, debemos de garantizar que aquel equilibrio que alcanzamos es el más deseable. De acuerdo con esta regla, el tipo de interés fijado por el banco central depende del output gap y la inflación porque se asume que la autoridad monetaria desea eliminar la fluctuación en ambas variables. 

WOODFORD (1999) argumentará que tanto la inflación como el output gap son objetivos deseables de la política monetaria. Para demostrarlo derivará una función de pérdida del banco central con la siguiente forma:\footnote{La derivará través de la aproximación de TAYLOR de segundo orden de la función de utilidad esperada del agente representativo de la economía lo que resulta fundamental ya que ésta no va a depender de funciones de pérdida \textit{ad hoc}, sino que se derivará de condiciones de optimización microfundamentadas del comportamiento de los hogares.
}

\eqblock{
\begin{aligned}
    E_0\left[\sum_{t=0}^{+\infty}\beta^t\;\overbrace{(\pi_t^2+\lambda\;\hat y^2)}^{L_t}\right]
\end{aligned}
}{Condición de deseabilidad del Equilibrio. Modelo NEK-2}

La razón de que la función de perdida de la autoridad monetaria depende de inflación y output gap es la siguiente: 
\begin{lnum}
    \item \textit{Inflación}. En los modelos de la NEK, la inflación es un indicador directo de la ineficiencia en la asignación de recursos en la economía. Bajo plena flexibilidad de precios, todas las empresas fijan un precio igual y los consumidores consumen una cantidad igual de cada variedad del bien. No obstante, cuando existe inflación debido a que los precios se fijan à la CALVO, algunas empresas no podrán adaptar su precio al nuevo nivel de precios de equilibrio. Ello generará una dispersión ineficiente de precios, que conducirá a variaciones en la cesta d ellos consumidores, que a su vez traerán pérdidas directas de utilidad. Por tanto, en estos modelos la inflación es un mal a evitar.
    \item \textit{Output gap}. Éste se define como la diferencia entre el PIB nominal y el PIB potencial. Debido a que el nivel de producción alcanzado bajo plena flexibilidad de precios (PIB potencial) constituye el óptimo de Pareto, tal y como ocurre en los modelos DSGE de la NMC, toda la desviación de éste será subóptima en términos de bienestar. Por tanto, el output gap también es una señal de las ineficiencias presentes en la economía.
\end{lnum}

Pero, ¿es posible alcanzar de manera simultánea la estabilización del output y la inflación?\footnote{Nótese que, en cierto modo, esta cuestión ha constituido el centro del debate macroeconómico de la segunda mitad del siglo XX y se remonta a las discusiones iniciales sobre la Curva de PHILLIPS, donde se interpretaba que existía un trade off entre inflación y desempleo.
} De acuerdo con el modelo de la NEK no existiría ningún trade off entre inflación y estabilización pero, bajo ciertas condiciones y en ausencia de imperfecciones reales relevantes, los bancos centrales pueden usar la política monetaria para estabilizar simultáneamente inflación y actividad, sin conflicto entre estos dos objetivos, fijando el tipo de interés de modo que replique la evolución del tipo de interés natural. Esto es lo que se ha conocido como la \textit{Divina Coincidencia} (hallada por ROTEMBERG y WOODFORD (2005), pese a que el término fue acuñado por BLANCHARD y GALÍ (2005)).


Esto supone que los bancos centrales tienen que centrarse únicamente en el objetivo de inflación y conseguirán simultáneamente estabilizar la actividad. 

En este sentido, la política monetaria no debería modificar el tipo de interés nominal hasta niveles en los que llegara a ser inconsistente con los acontecimientos reales de la economía\footnote{De esta manera, la curva IS ampliada por las expectativas pasa a determinar un output gap presente nulo, lo que a su vez genera, de acuerdo con la curva de Phillips de la NEK una inflación presente nula. El resultado de aplicar una política monetaria óptima es que la economía obtendrá en cada período los mismos resultados en términos de output, empleo, ahorro y el resto de variables macroeconómicas relevantes, que en un contexto de plena flexibilidad de precios, maximizando así el bienestar de los agentes y alcanzando las asignaciones óptimas de Pareto.
}\textsuperscript{,}\footnote{Esta visión nos ofrece un marco para analizar el tono que está teniendo la política monetaria:
\begin{la}
    \item Si el tipo de interés nominal fijado por el banco central menos las expectativas de inflación es menor que el tipo de interés natural ($i_t-E[\pi_{t+1\;|\;\Omega_t}< r_t^n]$), la política monetaria tendrá un carácter expansivo, lo que generará un output gap positivo, debido a que los agentes percibirán estos tipos como un desincentivo al ahorro, aumentando su demanda agregada presente, que derivará en presiones inflacionistas contemporáneas.
    \item Si el tipo de interés nominal fijado por el banco central menos las expectativas de inflación es mayor que el tipo de interés natural ($i_t-E[\pi_{t+1\;|\;\Omega_t}> r_t^n]$), la política monetaria tendrá un carácter contractivo, lo que producirá un output gap negativo debido a que los agentes percibirán un encarecimiento del coste de oportunidad del consumo presente, que se traducirá en una caída de la demanda agregada y en presiones deflacionistas.
\end{la}
}. Como hay cierto grado de rigidez en los precios pero el banco central fija su tipo de interés para que la inflación sea cero, entonces la rigidez de precios no causa daño en agregado.

En un mundo de inflación cero, todavía algunas empresas desearían aumentar sus precios, mientras que otras desearían reducirlo, y el ajuste costoso en precios implicaría que algunas empresas terminarían produciendo más de lo deseado y otras menos de lo deseado. 

Sin embargo, el output gap agregado terminaría siendo cero, minimizando el efecto de la fricción. 

En este sentido, el modelo neokeynesiano podría denominarse neowickselliano, debido a que fue WICKSELL el primero en proponer que la política monetaria debería ser tal que el tipo de interés real replicará al tipo de interés natural de la economía, ya que en caso contrario se abriría el output gap y la economía se recalentaría (si el tipo de interés real es inferior al tipo de interés natural, generando inflación y aumento de la actividad) o se enfriaría (si el tipo de interés real es superior al tipo de interés natural, generando deflación y reducción de la actividad). 

Esta implicación de la inexistencia de un trade off entre ambas variables es poco probable en la realidad, de modo que se han propuesto distintas soluciones o modificaciones al modelo canónico:\footnote{However, the previous extreme result holds only when the flexible price (or natural) equilibrium allocation is optimal –that is, when nominal rigidities are the only distortion in the economy. More generally, the presence of real frictions is likely to drive a wedge between the natural and efficient levels of output. As a result, the steady state itself may be inefficient, or the presence of teal frictions may imply an inefficient response of natural output to some shocks, or both. As a result, a trade-off emerges between price stability and the attainment of an efficient level of economic activity, thus giving rise to a nontrivial optimal monetary policy problem.
} 
\begin{la}
    \item Una posibilidad, propuesta por MANKIW, REIS y BALL, es introducir una perturbación externa ($e_t$) en la curva de PHILLIPS de la NEK y suponer que pueden ser persistentes y presentar correlación serial, que puede ser interpretada como variaciones exógenas en impuestos distorsionantes o cambios exógenos en los salarios deseados o en los mark-ups de precios (en la CPNEK presetada ya contamos con este término: $\pi_t=\beta\;E_t[\pi_{t+1}\;|\;\Omega_t]+\delta\;x_t+e_t$)
    \item Otra solución fue propuesta por BLANCHARD y GALÍ (2005), quienes llegan a la conclusión de que, si el modelo canónico se amplía para introducir rigideces reales en los salarios,\footnote{BLANCHARD y GALÍ (2005) asumen rigideces reales basadas en fricciones en la búsqueda y emparejamiento en el mercado de trabajo y para ello recurren al modelo de DIAMOND, MORTENSEN y PISSARIDES [\authorfont{Ver Tema 3.A.26}].} los bancos centrales sí tienen que hacer frente a un trade off entre la inflación y el output gap.\footnote{De hecho, según los autores este modelo ampliado tiene implicaciones más realistas que el modelo normal.\\
Entre las imperfecciones reales que pueden generar estos trade-offs, los autores se han centrado en el ajuste lento de los salarios reales y, en particular, en el hecho de que, si el salario responde en una relación inferior a uno ante cambios en la tasa marginal de sustitución entre consumo y ocio ante un shock de oferta, el nivel natural de output generará excesivas fluctuaciones en relación a su nivel de eficiencia.\\
De hecho, estabilizar completamente la inflación requeriría cerrar el gap entre el output y su nivel natural, lo que generaría fluctuaciones que reducirían el bienestar. En este contexto, una política monetaria acomodaticia sería lo más adecuado en el corto plazo, es decir, permitir que la inflación fuera algo superior que la que prevalecería bajo flexibilidad salarial. De lo contrario, el advenimiento de shocks adversos daría lugar a una reducción del output ineficientemente elevada.
}
\end{la}

\paragraph{Costes de la inflación y beneficios de la estabilidad de precios}
La NEK establece un marco riguroso para justificar las políticas destinadas a la estabilidad de precios. 

En primer lugar, la inflación es un indicador de la ineficiencia en el nivel de actividad por dos motivos:
\begin{lnum}
    \item Se origina en una desviación del output respecto a su nivel natural, desviación causada por la existencia de rigideces nominales. Por ello, incluso si el banco central no se preocupa de la inflación en sí misma, encontrará deseable estabilizarla de manera indirecta como una forma de cerrar el output gap; y,
    \item La inflación genera una ineficiente distribución de los recursos entre las empresas ya que, si hay una inflación positiva, las empresas tendrán que aumentar los precios en cada período, pero como no todas las empresas pueden cambiarlos por la rigidez de precios, los precios relativos variarán sin que esto esté justificado por un shock sectorial, generando cantidades subóptimas de los bienes producidos y consumidos.
\end{lnum} 

En este sentido, se justifica el argumento de inflación cero bajo cualquier circunstancia, independientemente de los costes que tuviese en términos de empleo o actividad. 

Sin embargo, también existen varios factores que justifican tener un nivel positivo de inflación:
\begin{lnum}
    \item El riesgo de generar un problema de zero-lower bound si los tipos nominales están muy bajos;
    \item Independientemente de los niveles deseados de inflación, si hay un shock relacionado con la inflación de costes, existirá un trade off entre la estabilización de precios y la del output gap.
\end{lnum} 

En este sentido, como la variación de ambas variables es una fuente de pérdida de bienestar, será óptimo para el banco central adoptar una política acomodaticia de la inflación, por lo que se suele recomendar fijar un objetivo positivo en el medio plazo.

\paragraph{El papel de las expectativas y las ganancias de la credibilidad}
En este tipo de modelos, las decisiones de fijación de precios y salarios se vuelven forward-looking, ya que los individuos reconocen que los precios y salarios fijados hoy permanecerán más allá de un período, afectando de esta forma tanto a la inflación como al output gap, que adoptan también una perspectiva forward-looking. Como consecuencia de ello, las decisiones de política económica que sean anticipadas tendrán influencia en los ingresos actuales y, por ello, el banco central puede beneficiarse de ser capaz de influir en las expectativas de los agentes.

Si se espera un shock positivo de inflación, el banco central podría aplicar una política restrictiva para reducir el output gap y así hacer frente al shock. Sin embargo, con este tipo de modelos, si el banco central se compromete de manera creíble a llevar a cabo acciones para luchar contra el shock esperado de inflación, puede reducir en menor medida el output gap en el presente, prometiendo que habrá menores output gaps en el futuro, y conseguir con ello el mismo resultado de control de la inflación. Por tanto, el output y la inflación reaccionan a cambios o anuncios de cambios futuros del instrumento de política monetaria. El banco central puede conseguir sus objetivos cambiando algo hoy mismo, o simplemente anunciando que lo hará en algún momento futuro, que de lugar al problema de credibilidad de los anuncios.

De acuerdo con los modelos de la NEK, un aumento de la credibilidad del banco central hace posible reducir simultáneamente la volatilidad del output gap y de la inflación. En otras palabras, si el banco central lleva a cabo una política de desinflación creíble, esto genera mejoras de bienestar.

En resumen, la comunicación del banco central sobre sus acciones futuras juega un papel esencial.

\paragraph{Importancia de los niveles de output y del tipo de interés como objetivos de política económica}
Los niveles naturales de output y de tipos de interés juegan un papel fundamental en el diseño de la política monetaria de la NEK. Desafortunadamente, la falta de observación de estas variables hace que sean difíciles de usar en la práctica; además, reemplazarlas por variables proxy puede provocar más daños que beneficios. Precisamente por ello, la importancia de estas variables naturales hace el desarrollo de los modelos DSGE particularmente útil ya que estos modelos pueden usarse para hacer inferencia de estas variables (aunque de modo impreciso).

\subsubsection{Política fiscal}
Respecto a la política fiscal, resulta muy ilustrativa la sentencia anticipada de SOLOW (2002), cuando afirmaba que la discusión seria sobre política fiscal había desaparecido del debate económico. 

En primer lugar, cabe destacar que en política fiscal, existen menos puntos en común entre la escuela de la NEK y la NMC que en política monetaria.

Además, el instrumental de análisis es el mismo tanto en política monetaria como en política fiscal: la curva de PHILLIPS de la NEK, la curva IS ampliada por las expectativas y la regla de tipos de interés. Por tanto, dado el carácter microfundamentado y la presencia de expectativas racionales, en el modelo canónico operan tres mecanismos de transmisión de la política fiscal a través del aumento del gasto público:
\begin{lnum}
    \item \textit{Efecto riqueza negativo}. Por una parte, los agentes anticipan los incrementos de impuestos futuros necesarios para financiar ese aumento del gasto público, produciéndose un efecto riqueza negativo, que lleva a los hogares a reducir su consumo y aumentar su oferta de trabajo;\footnote{
    Este efecto es el mismo que en los modelos neoclásicos y, por tanto, se verifica la Equivalencia Ricardiana.\\
    Los modelos de la NEK-2 aceptan la Equivalencia Ricardiana, lo que limita el efecto expansivo de la política fiscal por el lado de la demanda, y el efecto de la política fiscal depende en gran medida del carácter acomodaticio de la política monetaria.
    }
    \item \textit{Efecto causado por la existencia de rigideces nominales}. Por otra parte, la existencia de rigideces nominales hace que ante el aumento de la demanda agregada las empresas reaccionen incrementando su producción y aumentando su demanda de trabajo (ajuste en cantidades y en precios). Es decir, existe un efecto positivo sobre el output gap y también un efecto positivo sobre la inflación;\footnote{En el modelo neokeynesiano canónico, la existencia de competencia monopolística implica que el salario real no queda determinado sólo por la productividad marginal, sino que depende también del margen, por lo que puede existir una divergencia entre la productividad marginal y el salario real ($\partial F/\partial L_t=\mu\;(W_t/P_t)$).  La presencia de rigideces nominales de precios implica que, ante un incremento del coste marginal, los precios no podrán ajustarse en la misma proporción, por lo que el margen se reducirá (será contracíclico) y, por tanto, ante expansiones de la demanda el margen se reduce, manteniéndose el equilibrio en el mercado de trabajo sin necesidad de que se reduzca el salario real.\\
    Así, ante una expansión del gasto público, las empresas elevarán su producción y su demanda de trabajo. El aumento de la demanda de trabajo podrá compensar la mayor oferta de trabajo generada por el efecto riqueza negativo, produciéndose elevaciones de los salarios reales que compensarán parcialmente la caída del consumo por el efecto riqueza negativo.
    } y,
    \item \textit{Efecto respuesta de la política monetaria}. La autoridad monetaria, siguiendo la regla de TAYLOR, responde a las presiones inflacionistas que generan mayores salarios reales y eleva los tipos de interés. La subida de tipos de interés hará que los hogares reaccionen sustituyendo consumo por ahorro (efecto \textit{crowding out}), lo que llevará a una disminución del efecto expansivo de la política fiscal.
\end{lnum}

De este modo, es posible que un aumento del gasto público dé lugar a un aumento del salario real y del consumo, mientras que el valor del multiplicador (y el efecto total de la política fiscal) dependerá de:
\begin{lnum}
    \item \textit{Grado de rigidez de precios} (correlación entre el mark up y la posición cíclica). Cuanto mayor sea la rigidez de precios, mayor prociclicidad del margen, mayor aumento de la demanda de trabajo por parte de las empresas, mayor aumento de la producción y mayor efectividad de la política fiscal;
    \item \textit{Carácter acomodaticio de la política monetaria}. Cuanto menor sea la sensibilidad del banco central a las desviaciones de la inflación y el output gap respecto de su objetivo, menos incrementará los tipos de interés y mayor será la efectividad de la política fiscal; y,
    \item \textit{Persistencia del shock de gasto público}. Cuanto más persistente sea el gasto mayor será la inflación esperada y mayor la respuesta del banco central.
\end{lnum}

Además, el análisis reciente de la política fiscal en el marco de los modelos de la NEK se ha centrado en estudiar las condiciones bajo las cuales el multiplicador puede ser elevado, con especial atención a la existencia de restricciones de liquidez. Así, GALÍ, LÓPEZ SALIDO y VALLÉS (2007) plantean que su existencia actúa como un mecanismo que puede amplificar los efectos expansivos de la política fiscal.\footnote{Estos autores consideran un modelo de la NEK estándar con competencia monopolística y rigideces nominales de precios e introducen la existencia de dos tipos de consumidores: ricardianos y no ricardianos (que no pueden ahorrar o endeudarse debido a la existencia de restricciones de liquidez).\\
La existencia de competencia monopolística y de rigideces de precios hace que el salario real pueda elevarse en respuesta a una perturbación del gasto público. Así, un aumento del gasto público que eleve los salarios y el empleo incrementa la renta salarial de los consumidores no ricardianos que, dada la restricción de liquidez, reaccionarán consumiendo más.\\
De este modo, la interacción de rigideces de precios, competencia monopolística y consumidores no ricardianos hace que el consumo responda positivamente a las perturbaciones expansivas del gasto público.
}

\subsection{Limitaciones}
Sin embargo, como ya hemos comentado, la crisis financiera de 2007-2009 hizo que los fallos que había en el pensamiento económico dominante se hicieran más claros y condujeron a una tremenda pérdida de confianza, ya que, en ese momento, no se contaba con herramientas suficientes para reaccionar a una crisis financiera global. El ex presidente del Banco Central Europeo, Jean-Claude Trichet, llegó a afirmar en un coloquio con amigos que: \textit{Los instrumentos proporcionados por la teoría económica dominante habían resultado ser del todo inútiles para preveer e implementar soluciones que permitieran salir de la debacle originada en estos mercados. Es más, se habían sentido abandonados por las recomendaciones económicas dominantes.}

A lo largo de la crisis económica, los gobiernos tuvieron que usar diversos instrumentos de política que, hasta hace poco, eran considerados inadecuados. Una parte muy relevante de las críticas provino de la escuela neokeynesiana. En febrero de 2010, el FMI publicó un documento titulado \textit{Replanteando la Política Macroeconómica} (BLANCHARD et al., 2010), en el cual se indicaba que varios de los planteamientos de política económica anteriores a la crisis habían tenido fallos relevantes o incluso no eran correctos, entre los cuales podemos destacar tres:
\begin{lnum}
    \item En primer lugar, con una inflación en torno al 2\% en los países desarrollados, los tipos de interés nominales a corto plazo se fijan a un nivel muy bajo. En este caso, si hay una necesidad de flexibilizar la política monetaria, como en la reciente crisis, el tipo de interés nominal se sitúa rápidamente en su \textit{zero lower bound} generando una trampa de liquidez;
    \item En segundo lugar, desde el momento en que la política monetaria y el Quantitative Easing alcanzan sus límites, se espera que la crisis se prolongue más de lo esperado. En estas circunstancias, el retardo interno de la política fiscal ya no sería un impedimento para aplicar una expansión del gasto público. Por otro lado, las políticas fiscales expansivas que se adoptaron para evitar un colapso de los mercados financieros produjeron un nuevo problema: el fuerte aumento tanto del déficit público como de la deuda. La falta de espacio fiscal impidió que las medidas expansivas se prolongasen en el tiempo: la política fiscal debía ser contracíclica en todo momento, para generar superávits en épocas de bonanza y poder utilizar posteriormente la política fiscal expansiva en épocas de crisis, sin deteriorar excesivamente las cuentas públicas; y,
    \item La tercera cuestión que quedó clara fue la falta de neutralidad de los mercados financieros y la necesidad de establecer regulación en términos macroeconómicos. La supresión de la regulación en materia financiera ofreció incentivos al arbitraje y a la creación de nuevos instrumentos financieros, lo que permitió a los agentes financieros evitar algunas reglas prudenciales y condujo a una excesiva asunción de riesgos y a un exceso de apalancamiento.\footnote{Para más información acerca de los aceleradores de la crisis, véase: \href{https://www.filmaffinity.com/es/film112844.html}{Inside job (2010)}\\
    Algunos modelos Agent-Based han probado que, de haber mantenido los tipos constantes que se estaban ofreciendo en 1997 (altos) o habiendo regulado (e impedido) el uso intensivo nuevos instrumentos financieros de alta volatilidad y riesgo como los \href{https://en.wikipedia.org/wiki/Mortgage-backed_security}{\textit{mortgage based securities}} generados por instituciones de crédito privado hubieran evitado el colapso del sistema financiero estadounidense}
\end{lnum}

En síntesis, para los modelos de la NEK (o Nueva Sintesis), siendo estos el paradigma reinante hasta el inicio de la crisis, los mercados financieros eran tratados como un velo, sin ningún papel central en los modelos macroeconómicos.\footnote{Bajo el supuesto de que el tipo de interés a corto plazo está relacionado con el precio de los activos a través del arbitraje y que los efectos reales de la política monetaria tienen lugar a través de los tipos de interés y los precios de los activos, la intermediación financiera no tendría ningún tipo de relevancia macroeconómica excepto por el canal del crédito de los bancos comerciales. Aspectos tales como el exceso de apalancamiento de los agentes o el exceso de exposición a un determinado mercado no deberían tenerse en cuenta a la hora de implementar la política monetaria.
} Además, la regulación financiera no se consideraba como una herramienta macroeconómica, sino que, en el ámbito financiero, la regulación debía centrarse en las instituciones bancarias con el objetivo de corregir errores que provienen de los problemas de información asimétrica. Los riesgos de los mercados financieros eran vistos como exógenos con respecto al comportamiento de cada una de las instituciones individuales. Además, los precios de los activos, las condiciones crediticias y la macroeconomía eran vistas como independientes del comportamiento de las empresas financieras.

Al ponerse de manifiesto estas cuestiones, los autores de la NEK han ido realizando distintas aportaciones. El nuevo \textit{state of the art} podría organizarse en torno a varias cuestiones que han saltado al centro del debate económico. Tal y como resalta BLANCHARD en la tercera conferencia del FMI, convocada con la cuestión de replantearse la política económica (2016), los nuevos elementos esenciales del debate económico son los problemas que afronta la política monetaria, el nuevo papel de la política fiscal y la regulación macroprudencial. 

\subsubsection{Política monetaria}
CÚRDIA y WOODFORD (2010) admiten un fracaso al no haber tenido en cuenta el sistema financiero a la hora de diseñar la política monetaria y, por eso, sugieren un nuevo diseño de los modelos de la NEK. Otros estudios se ocupan de la no neutralidad del sistema financiero, la mayoría de ellos incorporando en los modelos DSGE de carácter neokeynesiano la modelización del crédito y la banca, lo que permite analizar su impacto en el bienestar macroeconómico y en la política monetaria.

\subsubsection{Política fiscal}
Cabe señalar que una vez que comenzó la crisis se hizo evidente que la misma era mucho más intensa que las anteriores. Se aplicaron rápidamente políticas monetarias expansivas, pero estas políticas se vieron limitadas con el problema del zero lower bound.

Consecuentemente, fue necesario recurrir a la política fiscal expansiva con el fin de fortalecer la demanda agregada. Sin embargo, como BLANCHARD (2008) reconoce, la estructura del modelo canónico de la NEK puede haber influido en la forma en que se diseñó la política fiscal: probablemente porque los bancos centrales son instituciones dotadas con grandes departamentos de investigación, se ha prestado mucha más importancia a los estudios sobre política monetaria que sobre política fiscal. En 2008, justo después de la quiebra de Lehman Brothers, el Departamento de Estudios del FMI publicó un artículo en el que se discutían las posibles formas en que la política fiscal debe aplicarse ante la situación de la crisis.

Sin embargo, como consecuencia de la crisis, el déficit público también se amplió excesivamente, al igual que la deuda pública, a lo que los mercados financieros globales han respondido de una manera muy adversa a este escenario, con el aumento de las primas de riesgo.\footnote{En este contexto de deterioro fiscal, algunos de los autores de la NEK-1, como TAYLOR, revivieron argumentos contra el activismo fiscal y se mostraron escépticos sobre la efectividad fiscal. Según TAYLOR, la política fiscal debe centrarse en la reducción del déficit y del crecimiento de la ratio deuda-PIB.
} No obstante, desde una perspectiva de medio plazo, BLANCHARD (2010) sostiene que, aunque los gobiernos deben demostrar un proyecto creíble para retirar el estímulo fiscal y políticas para volver a una senda de deuda pública sostenible, es necesario aumentar el crecimiento del gasto público en el corto plazo. Eso significa que, aun habiendo reconocido la necesidad de ajustar las cuentas públicas a medio plazo, BLANCHARD admite que una contracción fiscal con el fin de lograr una consolidación de la deuda pública es prematura y que la política fiscal ha contribuido de manera decisiva a evitar una caída de la producción más grave.

La conclusión más clara es que la política fiscal ha retornado al centro del debate económico:
\begin{lnum}
    \item \textit{Desde el punto de vista empírico}. Se cuestiona la validez y la aplicabilidad de consolidaciones fiscales expansivas, poniendo en entredicho las soluciones de austeridad fiscal (PEROTTI, 2011), y se estimaron multiplicadores fiscales de una manera innovadora, teniendo en cuenta los cambios normativos en los impuestos en Estados Unidos con el fin de evaluar con mayor precisión si esos cambios fueron promulgados con fines contracíclicos (ROMER y ROMER, 2010); y,
    \item \textit{Desde el punto de vista teórico}. La NEK se ha centrado en analizar el efecto de la política fiscal en las situaciones en las que la política monetaria está limitada, como cuando los tipos de interés nominales alcanzan el límite inferior (zero lower bound)\footnote{Véase: \href{https://www.frbsf.org/economic-research/wp-content/uploads/sites/4/2-Wu.pdf}{\textit{The Four Equation New Keynesian Model}. SIMS y WU, 2019}\\
    Además, se han estudiado las distintas respuestas de política monetaria a una política fiscal expansiva que pueden generar que el multiplicador del gasto público sea mayor que uno (WOODFORD, 2011), y analizado estos efectos de la política fiscal en una situación en la que se ha alcanzado el límite inferior de los tipos de interés, encontrando que el multiplicador fiscal puede tomar valores muy superiores a uno [\authorfont{Ver tema 3.A.38}] (CHRISTIANO et al., 2011)}. 
\end{lnum}

\subsubsection{Regulación macroprudencial}
Algunos autores de la NEK con un enfoque más práctico (BLANCHARD, BERNANKE y BLINDER) se han preocupado recientemente por la justificación de la regulación macroprudencial, que en su opinión debe ser coherente con el grado de sofisticación del sistema financiero y con el impacto macroeconómico que puede ejercer.\footnote{Entre los economistas académicos, es posible encontrar algunas contribuciones que revelan cierta preocupación respecto a la no neutralidad del sistema financiero: (i) GERTLER et al. (2010) dan un nuevo enfoque típicamente keynesiano a la regulación microprudencial; (ii) también incluyen en sus modelos la posibilidad de la política macroprudencial para contrarrestar el incentivo para la excesiva asunción de riesgos por parte de los bancos; y, (iii) una conclusión importante es que la regulación macroprudencial aumenta el bienestar de la sociedad.
} Además, en la mayor parte de los países se han desarrollado herramientas para medir el riesgo sistemático.

En realidad, las discusiones teóricas posteriores a la crisis se han alejado en cierta manera de los planteamientos de la NEK.

\clearpage
\section{Conclusión}
A modo de conclusión, y una vez vistas las aportaciones realizadas en algunas de las cuestiones más relevantes dentro del debate de política económica en la actualidad cabe plantearse si, para los autores de la NEK, sus planteamientos y sus modelos siguen siendo correctos. La respuesta sería que sí. 

En la mayoría de las publicaciones, explícita o implícitamente, los nuevos keynesianos manifiestan que los errores se pueden solventar con algunos cambios en los modelos. A pesar de la inadecuación a priori de los modelos DSGE para hacer frente a la política fiscal (como se indica por BLANCHARD, 2008), muchos macroeconomistas los están utilizando para acercarse a la eficacia y conveniencia de la política fiscal. Esto significa que se sigue considerando que este tipo de modelos siguen siendo adecuados para guiar la política macroeconómica. 

Sin embargo, no están exentos de críticas y algunos economistas señalan que se les ha otorgado excesiva relevancia. Pese a la crisis económica, la NEK sigue siendo una de las escuelas centrales, quizás la más relevante en el debate macroeconómico. Los modelos DSGE siguen usándose tanto en Europa como en Estados Unidos, las fricciones siguen teniendo un papel preeminente en los modelos y se ha conseguido introducir nuevos elementos de política fiscal y de los mercados financieros en el estudio.

\clearpage
\pagenumbering{gobble}
\MyBibliography
\MyBibItem{Autor}{Título}{Año}{DOI -- LINK}


% --- Anexos (no aparecen en el índice) ---
\clearpage
\anexo{\textbf{Preguntas Test}}
%\input{\TemaCode/questions.tex} 


\end{document}